\subsection{乘法自同态}

设 X 是局部紧拓扑空间，\(\mu\) 是 X 上的正测度。设 \(p \in [1, +\infty[\) 。

设 \(g \in \mathcal{L}^{\infty}(\mathrm{X},\mu)\) 。用 \(m_{g}\) 表示由 \(f \mapsto g \cdot f\) 所给出的从 \(\mathcal{L}^{p}(\mathrm{X},\mu)\) 到自身的映射。映射 \(m_{g}\) 是线性的且连续的，因为

\[
\mathrm{N} _ {p} (m _ {g} (f)) \leqslant \mathrm{N} _ {\infty} (g) \mathrm{N} _ {p} (f)\tag{1}
\]

对每个函数 \(f \in \mathcal{L}^{p}(\mathrm{X}, \mu)\) 成立。特别地，函数 \(m_{g}(f)\) 是 \(\mu\)-零测的，如果 f 是 \(\mu\)-零测的。于是，映射 \(m_{g}\) 通过过渡到商诱导出 \(\mathrm{L}^{p}(\mathrm{X}, \mu)\) 的一个自同态，我们将它记作 \(\widetilde{m}_{g}\) 。

引理 7. --- 设 \(g \in \mathcal{L}^{\infty}(\mathrm{X},\mu)\) 。自同态 \(\tilde{m}_g\) （\(\mathrm{L}^{p}(\mathrm{X},\mu)\) 的）是单射的，当且仅当由 \(x \in \mathrm{X}\) （满足 \(g(x) = 0\) ）所成的集是局部 \(\mu\)-零测的。

用 A 表示由 \(x \in X\) （满足 \(g(x) = 0\) ）所成的集。

假设 A 是局部 \(\mu\)-零测的。设 \(f \in \mathcal{L}^p(\mathrm{X}, \mu)\) ，\(\widetilde{f}\) 是它在 \(\mathrm{L}^p(\mathrm{X}, \mu)\) 中的类。假设 \(\widetilde{m}_g(\widetilde{f}) = 0\) 。按定义，这等价于函数 \(fg\) 是 \(\mu\)-零测的，于是由 \(x \in \mathrm{X}\) （满足 \(f(x)g(x) \neq 0\) ）所成的集 B 是 \(\mu\)-零测的。函数 \(f\) 在 \(\mathrm{A} \cup \mathrm{B}\) 之外为零，故是 \(\mu\)-零测的。这就蕴含自同态 \(\widetilde{m}_g\) 是单射的。

反之，假设 A 不是局部 \(\mu\)-零测的。设 C 是 X 的紧子集，使得 A∩C 不是 \(\mu\)-零测的，并设 \(\varphi\) 是 A∩C 的特征函数。函数 \(\varphi\) 在 L \(^{p}\) (X, \(\mu\) ) 中的类不为零，但 \(m_{g}(\varphi)\) 的类为零，故 \(m_{g}\) 不是单射的。

由于局部 \(\mu\)-零测函数与 \(\mu\)-零测函数之积是 \(\mu\)-零测的，自同态 \(\widetilde{m}_g\) 只依赖于 \(g\) 在 \(\mathrm{L}^{\infty}(\mathrm{X},\mu)\) 中的类。对 \(\widetilde{g}\in \mathrm{L}^{\infty}(\mathrm{X},\mu)\) ，我们也用 \(\widetilde{m}_{\widetilde{g}}\) 记自同态 \(\widetilde{m}_g\) （\(\mathrm{L}^p (\mathrm{X},\mu)\) 的），这里 \(g\in \mathcal{L}^{\infty}(\mathrm{X},\mu)\) 是以 \(\widetilde{g}\) 为类的任一函数。我们称 \(\widetilde{m}_{\widetilde{g}}\) 为由 \(\widetilde{g}\) 所作的、\(\mathrm{L}^p (\mathrm{X},\mu)\) 中的乘法自同态。

命题 5. --- 设 \(p \in [1, +\infty[\) 。映射 \(m: g \mapsto \widetilde{m}_g\) （从 \(\mathrm{L}^{\infty}(\mathrm{X}, \mu)\) 到 \(\mathcal{L}(\mathrm{L}^{p}(\mathrm{X}, \mu))\) 中的）是含幺巴拿赫代数的一个等距态射。

对 \(a\) ， \(b\) （属于 \(\mathbf{C}\) ）以及 \(g_{1}\) ， \(g_{2}\) （属于 \(\mathrm{L}^{\infty}(\mathrm{X},\mu)\) ），可直接验证 \(\widetilde{m}_{ag_1+bg_2}=a\widetilde{m}_{g_1}+b\widetilde{m}_{g_2}\) ，故映射 \(\widetilde{m}\) 是线性的。此外，我们有 \(\widetilde{m}_{1}=1_{\mathrm{L}^{p}(\mathrm{X},\mu)}\) 与 \(\widetilde{m}_{g_{1}g_{2}}=\widetilde{m}_{g_{1}}\widetilde{m}_{g_{2}}\) ，故映射 \(\widetilde{m}\) 是从 \(\mathrm{L}^{\infty}(\mathrm{X},\mu)\) 到 \(\mathcal{L}(\mathrm{L}^{p}(\mathrm{X},\mu))\) 中的含幺代数态射。

上面的公式 (1) 表明 \(\widetilde{m}\) 的范数 \(\leqslant 1\) 。

设 \(g \in \mathcal{L}^{\infty}(\mathrm{X},\mu)\) ，\(\widetilde{g}\) 是它在 \(\mathrm{L}^{\infty}(\mathrm{X},\mu)\) 中的类。设 \(\varepsilon > 0\) 。由 \(x \in X\) （满足 \(|g(x)| > \| \widetilde{g}\|_{\infty} - \varepsilon\) ）所成的集 Y 不是局部 \(\mu\)-零测的。因此存在 X 的紧子集 C 使得 \(\mu (\mathrm{Y} \cap \mathrm{C}) > 0\) 。设 \(\varphi\) 是 \(\mathrm{Y} \cap \mathrm{C}\) 的特征函数。由于

\[
\begin{array}{c} \| \widetilde {m} _ {\widetilde {g}} (\varphi) \| _ {p} = \Bigl (\int_ {\mathrm{X}} | \varphi g | ^ {p}   d \mu \Bigr) ^ {1 / p} \geqslant \Bigl (\int_ {\mathrm{Y}} | \varphi g | ^ {p}   d \mu \Bigr) ^ {1 / p} \\ \geqslant (\| \widetilde {g} \| _ {\infty} - \varepsilon) \mu (\mathrm{Y} \cap \mathrm{C}) ^ {1 / p} = (\| \widetilde {g} \| _ {\infty} - \varepsilon) \| \varphi \| _ {p}, \end{array}
\]

便得到 \(\|\widetilde{m}_{\widetilde{g}}\|\geqslant\|\widetilde{g}\|_{\infty}-\varepsilon\) 。由于 \(\varepsilon\) 是任意的，可推出 \(\|\widetilde{m}_{\widetilde{g}}\|\geqslant\|\widetilde{g}\|_{\infty}\) ，这就证明了 \(\widetilde{m}\) 是等距的。

引理 8. --- 设 \((g_n)\) 是 \(\mathcal{L}^\infty(\mathrm{X},\mu)\) 中 \(\mu\)-几乎处处逐点收敛的有界序列。设 \(\widetilde{m} \in \mathrm{L}^\infty(\mathrm{X},\mu)\) 是其极限的类。那么 \(\widetilde{m}_{g_n}\) 收敛于 \(\widetilde{m}_{\widetilde{g}}\) ，这里是在赋有逐点收敛拓扑的空间 \(\mathcal{L}(\mathrm{L}^p(\mathrm{X},\mu))\) 中。

设 \(f \in \mathcal{L}^p(\mathrm{X},\mu)\) 。序列 \((g_{n}f)\) 在 \(\mathcal{L}^p(\mathrm{X},\mu)\) 中有界，且 \(\mu\)-几乎处处逐点收敛于 \(gf\) 。由勒贝格定理（INT, IV, p. 137, § 3, no. 7，定理 6），序列 \((g_{n}f)\) 收敛于 \(gf\) （在 \(\mathcal{L}^p(\mathrm{X},\mu)\) 中），断言得证。

现在考虑 p = 2 的情形。

命题 6. --- 映射 m: \(g \mapsto \widetilde{m}_{g}\) 是从星代数 \(\mathrm{L}^{\infty}(\mathrm{X}, \mu)\) 到星代数 \(\mathcal{L}(\mathrm{L}^{2}(\mathrm{X}, \mu))\) 中的含幺等距态射。

特别地，对每个 \(g \in \mathrm{L}^{\infty}(\mathrm{X}, \mu)\) ，乘法自同态 \(\widetilde{m}_{g}\) 是正规的；它是埃尔米特的，当且仅当 g \(\mu\)-几乎处处局部取实值。

由命题 5，映射 \(\tilde{m}\) 是从 \(\mathrm{L}^{\infty}(\mathrm{X},\mu)\) 到 \(\mathcal{L}(\mathrm{L}^{2}(\mathrm{X},\mu))\) 中的含幺巴拿赫代数的单射等距态射。设 \(g\in \mathcal{L}^{\infty}(\mathrm{X},\mu)\) 。对 \(f_{1}\) 与 \(f_{2}\in \mathcal{L}^{2}(\mathrm{X},\mu)\) ，我们有

\[
\langle f _ {1} \mid \widetilde {m} _ {g} (f _ {2}) \rangle = \int_ {\mathrm{X}} \overline {{f _ {1} (x)}} g (x) f _ {2} (x) d \mu (x) = \langle \widetilde {m} _ {\overline {{g}}} (f _ {1}) \mid f _ {2} \rangle ,
\]

由此得出 \(\widetilde{m}_g^* = \widetilde{m}_{\overline{g}}\) ，这就证明了 \(m\) 是对合态射。最后的断言由此得出（参见 I, p. 106，命题 5）。

推论. --- 设 \(g \in \mathrm{L}^{\infty}(\mathrm{X},\mu)\) 。

\begin{enumerate}
\def\labelenumi{\alph{enumi})}
\item
  \(\widetilde{m}_g\) 在 \(\mathcal{L}(\mathrm{L}^2 (\mathrm{X},\mu))\) 中的谱是 \(\mu\)-本质意义下 \(g\) 的像；
\item
  自同态 \(\widetilde{m}_g\) 是正的，当且仅当 \(g\) \(\mu\)-几乎处处局部为正；
\item
  对任意函数 \(f \in \mathcal{C}(\mathrm{Sp}(\widetilde{m}_g))\) ，我们有 \(f(\widetilde{m}_g) = \widetilde{m}_{f \circ g}\) 。
\end{enumerate}

由于 \(\widetilde{m}\) 是单射的，由 I, p. 106 的命题 5，我们有 \(\mathrm{Sp}(\widetilde{m}_{g}) = \mathrm{Sp}(g)\) ；结论由 IV, p. 185 的引理 6、IV, p. 185 的命题 4 以及 I, p. 115 的定义 6 得出。

命题 7. --- 态射 \(\tilde{m}\) （从 \(\mathrm{L}^{\infty}(\mathrm{X},\mu)\) 到 \(\mathcal{L}(\mathrm{L}^{2}(\mathrm{X},\mu))\) 中的）的像是代数 \(\mathcal{L}(\mathrm{L}^{2}(\mathrm{X},\mu))\) 的一个极大交换子代数。

只须证明：若 \(u \in \mathcal{L}(\mathrm{L}^{2}(\mathrm{X}, \mu))\) 是与 \(\widetilde{m}_{g}\) （对每个函数 \(g \in \mathcal{L}^{\infty}(\mathrm{X}, \mu)\) ）都交换的自同态，则 u 属于 \(\widetilde{m}\) 的像。设 u 是这样一个自同态。

我们用 \(\tilde{f}\) 表示 \(\mathrm{L}^2 (\mathrm{X},\mu)\) 中函数 \(f\in \mathcal{L}^2 (\mathrm{X},\mu)\) 的类。用 \(v\) 表示从 \(\mathcal{L}^2 (\mathrm{X},\mu)\) 到 \(\mathrm{L}^2 (\mathrm{X},\mu)\) 中的、由 \(f\mapsto u(\widetilde{f})\) 所定义的线性映射。我们有 \(v\circ m_g = \widetilde{m}_g\circ v\) （对所有 \(g\in \mathcal{L}^{\infty}(\mathrm{X},\mu)\) ）。

对 X 的任意 \(\mu\)-可测子集 Y，我们用 \(\varphi_{\mathrm{Y}}\) 表示它的特征函数；自同态 \(\widetilde{m}_{\varphi_{\mathrm{Y}}}\) 是 \(\mathrm{L}^2 (\mathrm{X},\mu)\) 的一个正交投影算子。

首先假设 X 是紧的，从而常值函数 1 的类属于 \(\mathrm{L}^{2}(\mathrm{X},\mu)\) 。设 g 是 \(\mathcal{L}^{2}(\mathrm{X},\mu)\) 中的函数，它在 \(\mathrm{L}^{2}(\mathrm{X},\mu)\) 中的类为 \(v(1)\) 。

设 \(c\) 是正实数，Y 是由 \(x \in \mathrm{X}\) （满足 \(|g(x)| \geqslant c\) ）所成的集；它是一个 \(\mu\)-可测集。我们在 \(\mathrm{L}^2(\mathrm{X}, \mu)\) 中有等式 \(\widetilde{\varphi_{\mathrm{Y}}} g = \widetilde{m}_{\varphi_{\mathrm{Y}}}(v(1)) = v(m_{\varphi_{\mathrm{Y}}}(1)) = v(\varphi_{\mathrm{Y}})\) 。

又由于 \(|c\varphi_{\mathrm{Y}}| \leqslant |g\varphi_{\mathrm{Y}}|\) ，我们得到

\[
c ^ {2} \mu (\mathrm{Y}) = \int_ {\mathrm{X}} (c \varphi_ {\mathrm{Y}}) ^ {2} d \mu \leqslant \int_ {\mathrm{X}} | \varphi_ {\mathrm{Y}} g | ^ {2} d \mu = \| \widetilde {u} (\varphi_ {\mathrm{Y}}) \| ^ {2} \leqslant \| u \| ^ {2} \mu (\mathrm{Y}).
\]

这个不等式表明 \(\mu(\mathrm{Y}) = 0\) （若 \(c > \|u\|\) ），从而函数 \(g\) \(\mu\)-几乎处处以 \(\|u\|\) 为界。最后，对每个函数 \(f \in \mathcal{C}(\mathrm{X})\) ，我们有

\[
u (\widetilde {f}) = v (m _ {f} (1)) = \widetilde {m} _ {f} (v (1)) = \widetilde {m} _ {f} (\widetilde {g}) = \widetilde {m} _ {g} (\widetilde {f}),
\]

由此得到等式 \(u = \tilde{m}_{g}\) ，特别地 \(\mathrm{N}_{\infty}(g) = \| u\|\) 。

现在考虑一般情形。存在 X 的两两不相交紧子集所成的局部可数族 \((\mathrm{X}_{i})_{i\in\mathrm{I}}\) ，使得 \(Z = X - \bigcup_{i \in I} X_i\) 是局部 \(\mu\)-零测的（INT, IV, p.190, § 5, no. 9，命题 14）。设 \(\mu_i\) 是 \(\mu\) 在 \(X_i\) 上诱导的测度（INT, IV, p. 186, § 5, no. 7，定义 4）。

由 IV, p. 179 的命题 1，对每个 \(i \in I\) ，像 \(\mathrm{E}_i\) （\(\widetilde{m}_{\varphi_{\mathrm{X}_i}}\) 的）与 \(\mathrm{L}^2(\mathrm{X}_i, \mu_i)\) 等同——通过 \(f \mapsto f|\mathrm{X}_i\) 。对每个函数 \(g \in \mathcal{L}^\infty(\mathrm{X}, \mu)\) ，乘法自同态 \(\widetilde{m}_g\) 与 \(\widetilde{m}_{\varphi_{\mathrm{X}_i}}\) 交换，因而使 \(\mathrm{E}_i\) 保持稳定，且 \(\mathrm{E}_i\) 的、由 \(\widetilde{m}_g\) 过渡到子空间而得的自同态与乘以 \(g|\mathrm{X}_i\) 的乘法自同态（\(\mathrm{L}^2(\mathrm{X}_i, \mu_i)\) 上的）重合。

由于 \(u\) 与 \(m_{\varphi_{X_i}}\) 交换，它也使子空间 \(\mathrm{E}_i\) 保持稳定。用 \(u_i\) 表示 \(\mathrm{L}^2 (\mathrm{X}_i,\mu_i)\) 的、由 \(u\) 过渡到子空间而得的自同态。

由于从 \(\mathcal{L}^{\infty}(\mathrm{X},\mu)\) 到 \(\mathcal{L}^{\infty}(\mathrm{X}_{i},\mu_{i})\) 中的、由 \(g\mapsto g|\mathrm{X}_{i}\) 所定义的映射是满射的，假设蕴含 \(u_{i}\) 与 \(\widetilde{m}_{g}\) 交换（对每个函数 \(g\in \mathcal{L}^{\infty}(\mathrm{X}_{i},\mu_{i})\) ）。由证明的第一部分，存在函数 \(g_{i}\in \mathcal{L}^{\infty}(\mathrm{X}_{i},\mu_{i})\) ，使得 \(u_{i} = \widetilde{m}_{g_{i}}\) 且 \(\mathrm{N}_{\infty}(g_{i})\leqslant \| u_{i}\| \leqslant \| u\|\) 。

X 上与 \(g_{i}\) 在 \(X_{i}\) 上重合、在 Z 上为零的函数 g 是有界且 \(\mu\)-可测的（INT, IV, p. 193, § 5, no. 10，命题 16），故 \(g \in \mathcal{L}^{\infty}(X, \mu)\) 。对每个 \(i \in I\) ，u 在 \(E_{i}\) 上的限制与 \(\tilde{m}_{g}\) 的限制重合；因此，由 IV, p. 179 的命题 1, c)，我们有 \(u = m_{g}\) 。

推论. --- 设 u 是希尔伯特空间 \(\mathrm{L}^{2}(\mathrm{X},\mu)\) 的自同态，它与 \(\widetilde{m}_{g}\) （对每个函数 \(g\in\mathcal{K}(\mathrm{X})\) ）交换。那么存在唯一的元素 \(f\in\mathrm{L}^{\infty}(\mathrm{X},\mu)\) ，使得 \(u=\widetilde{m}_{f}\) 。

由命题 7，只须证明 \(u\) 与 \(\widetilde{m}_{g}\) 交换（对所有 \(g\in\mathcal{L}^{\infty}(\mathrm{X},\mu)\) ）。设 \(h_{1}\) 与 \(h_{2}\) 是 \(\mathcal{L}^{2}(\mathrm{X},\mu)\) 中的元素，它们的类 \(\widetilde{h}_{1}\) 与 \(\widetilde{h}_{2}\) 属于 \(\mathrm{L}^{2}(\mathrm{X},\mu)\) 。设 \(k_{1}\) （相应地， \(k_{2}\) ）是 \(\mathcal{L}^{2}(\mathrm{X},\mu)\) 中的函数，其类为 \(u(\widetilde{h}_{1})\) （相应地， \(u^{*}(\widetilde{h}_{2})\) ）。

设 \(h = h_1 \overline{k}_2 - k_1 \overline{h}_2\) ；我们有 \(h \in \mathcal{L}^1(\mathrm{X}, \mu)\) 。定义测度 \(\nu = h \cdot \mu\) （在 \(\mathrm{X}\) 上）；它是有界的。对每个 \(g \in \mathcal{L}^\infty(\mathrm{X}, \mu)\) ，我们有

\[
\begin{array}{c} \langle \widetilde {h} _ {2} | u (\widetilde {m} _ {g} (\widetilde {h} _ {1})) - \widetilde {m} _ {g} (u (\widetilde {h} _ {1})) \rangle = \langle u ^ {*} (\widetilde {h} _ {2}) | \widetilde {m} _ {g} (\widetilde {h} _ {1}) \rangle - \langle \widetilde {h} _ {2} | \widetilde {m} _ {g} (u (\widetilde {h} _ {1})) \rangle \\ = \int_ {\mathrm{X}} g \cdot h _ {1} \cdot \overline {{k}} _ {2} d \mu - \int_ {\mathrm{X}} g \cdot k _ {1} \cdot \overline {{h}} _ {2} d \mu = \nu (g). \end{array}
\]

于是，按假设有 \(\nu(g) = 0\) （对每个函数 \(g \in \mathcal{K}(\mathrm{X})\) ），即 \(\nu = 0\) 。由此可得

\[
\langle \widetilde {h} _ {2} | u (\widetilde {m} _ {g} (\widetilde {h} _ {1})) - \widetilde {m} _ {g} (u (\widetilde {h} _ {1})) \rangle = \nu (g) = 0
\]

对所有 \(g \in \mathcal{L}^{\infty}(\mathrm{X},\mu)\) 成立。由于这对所有 \(h_1\) 与 \(h_2\) （\(\mathcal{L}^2 (\mathrm{X},\mu)\) 中的）都成立，我们有 \(u\circ \widetilde{m}_g = \widetilde{m}_g\circ u\) 。

注. --- 在下文中，我们常径直用 \(m_{g}\) 表示乘以 g 的乘法算子（\(\mathrm{L}^{p}(\mathrm{X},\mu)\) 上的）。

\subsection{谱测度}

在本节中，我们固定一个复希尔伯特空间 E。

回顾：若 A 是交换含幺星代数，我们用 X(A) 表示其幺特征标构成的紧拓扑空间（I, p. 29 的推论 1），并用 \(\mathcal{G}_{\mathrm{A}}\) 表示 A 的盖尔范德变换（I, p. 7 的定义 5），它是从 A 到 \(\mathcal{C}(\mathrm{X}(\mathrm{A}))\) 上的含幺对合代数的等距同构（I, p. 108 的定理 1）。我们将用 \(\mathcal{H}_{\mathrm{A}}\) 表示其逆同构。

引理 9. --- 设 A 是 \(\mathcal{L}(\mathrm{E})\) 的交换含幺星子代数。对 E 中所有 \(x\) 与 \(y\) ，映射 \(\mu\) ——从 \(\mathsf{X}(\mathrm{A})\) 到 \(\mathbf{C}\) 中，由 \(\mu(f) = \langle x | \mathcal{H}_{\mathrm{A}}(f)y \rangle\) （对每个 \(f \in \mathcal{C}(\mathsf{X}(\mathrm{A}))\) ）所定义——是紧空间 \(\mathsf{X}(\mathrm{A})\) 上的有界测度。当 \(x = y\) 时它是正的。其范数 \(\leqslant \|x\|\ \|y\|\) ，且当 \(x = y\) 时取等号。

映射 \(\mu\) 是线性的。对每个函数 \(f\in\mathcal{C}(\mathsf{X}(\mathsf{A}))\) ，我们有

\[
| \mu (f) | \leqslant \| x \| \| y \| \| \mathcal {H} _ {\mathrm{A}} (f) \| = \| x \| \| y \| \| f \|,
\]

因而 \(\mu\) 是 \(\mathsf{X}(\mathsf{A})\) 上范数 \(\leqslant \| x\|\| y\|\) 的有界测度。

假设 \(x = y\) 。对每个正函数 \(f \in \mathcal{C}(\mathsf{X}(\mathsf{A}))\) ，元素 \(\mathcal{H}_{\mathsf{A}}(f)\) 是 A 的正元素（I, p. 115 的例 1），因而是 \(\mathcal{L}(\mathsf{E})\) 的正元素，于是 \(\mu(f) = \langle x | \mathcal{H}_{\mathsf{A}}(f)(x) \rangle \geqslant 0\) （I, p. 138 的命题 8）。这表明 \(\mu\) 是正测度。由于 \(\mu(1) = \|x\|^2\) ，\(\mu\) 的总质量是 \(\|x\|^2\) （INT, III, p. 58, § 1, no. 8，推论 2）。

定义 2. --- 设 A 是 \(\mathcal{L}(\mathrm{E})\) 的交换含幺星子代数。对 E 中所有 \(x\) 与 \(y\) ，线性形式 \(f \mapsto \langle x | \mathcal{H}_{\mathrm{A}}(f)y \rangle\) （\(\mathcal{C}(\mathrm{X}(\mathrm{A}))\) 上的）称为 \((x,y)\) 关于 A 的谱测度。若 \(x = y\) ，我们称它为 \(x\) 关于 A 的谱测度。

注. --- 设 A 是 \(\mathcal{L}(\mathrm{E})\) 的交换含幺星子代数。对 E 中所有 \(x\) 与 \(y\) ，用 \(\mu_{x,y}\) 表示 \((x,y)\) 关于 A 的谱测度。由 \((x,y) \mapsto \mu_{x,y}\) 定义的从 \(\mathrm{E} \times \mathrm{E}\) 到 \(\mathcal{M}^1(\mathrm{X}(\mathrm{A}))\) 中的映射是半双线性的。

设 \(u\) 是 E 的正规自同态，A 是 \(\mathcal{L}(\mathrm{E})\) 的、由 \(u\) 所生成的含幺星子代数。它是交换的。空间 \(\mathsf{X}(\mathrm{A})\) 与 \(\mathrm{Sp}(u)\) 等同——通过映射 \(\chi \mapsto \chi(u)\) （I, p. 109 的引理 10）。于是，谱测度 \(\mu_{x,y}\) ——\((x,y)\) 关于 A 的——与 \(\mathrm{Sp}(u)\) 上的一个测度等同，该测度称为 \((x,y)\) 关于 \(u\) 的谱测度。于是，对每个函数 \(f \in \mathcal{C}(\mathrm{Sp}(u))\) ，我们有

\[
\int_ {\mathrm{Sp} (u)} f d \mu_ {x, y} = \langle x | f (u) y \rangle
\]

依据 I, p. 110 的定义 5。

\subsection{希尔伯特空间上自同态的交换星代数}

定义 3. --- 设 A 是 \(\mathbf{C}\) 上的代数，E 是复拓扑向量空间。设 π 是 A 在 E 中的一个表示。若对 a ∈ A 的元素 π(a)x 所成的集在 E 中是全的，则称 E 的元素 x 是 π 的循环向量。

设 \(u \in \mathcal{L}(\mathrm{E})\) 是 E 的自同态。我们称 \(x \in \mathrm{E}\) 是 \(u\) 的循环向量，如果它是由 \(u\) 在 \(\mathcal{L}(\mathrm{E})\) 中所生成的星子代数的恒等表示的循环向量。

命题 8. --- 设 E 是复希尔伯特空间。设 A 是 \(\mathcal{L}(\mathrm{E})\) 的交换含幺星子代数。设 \(x\) 是 E 的元素。令 \(\mathrm{E}_x = \overline{\mathrm{A} \cdot x}\) ，并用 \(\mu_x\) 表示 \(x\) 关于 A 的谱测度。

存在唯一的等距同构 \(\theta_x\) ——从 \(\mathrm{L}^2 (\mathsf{X}(\mathrm{A}),\mu_x)\) 到 \(\mathrm{E}_x\) 上——使得 \(\theta_{x}(f) = \mathcal{H}_{\mathrm{A}}(f)(x)\) 对每个函数 \(f\in \mathcal{C}(\mathsf{X}(\mathrm{A}))\) 成立。对每个元素 \(a\in \mathrm{A}\) ，空间 \(\mathrm{E}_x\) 在 \(a\) 下稳定，且若用 \(a_{x}\) 表示由 a 过渡到子空间而得的 \(\mathrm{E}_x\) 的自同态，则我们有 \(a_{x}\circ \theta_{x} = \theta_{x}\circ m_{\mathcal{G}_{\mathrm{A}}(a)}\) 。

设 \(\widetilde{\theta}_x\) 是从 \(\mathcal{C}(\mathrm{X(A)})\) 到 \(\mathrm{E}_x\) 中的线性映射，定义为

\[
\widetilde {\theta} _ {x} (f) = \mathcal {H} _ {\mathrm{A}} (f) (x).
\]

对每个函数 \(f \in \mathcal{C}(\mathsf{X}(\mathsf{A}))\) ，我们有

\[
\| \widetilde {\theta} _ {x} (f) \| ^ {2} = \langle \mathcal {H} _ {\mathrm{A}} (f) x | \mathcal {H} _ {\mathrm{A}} (f) x \rangle = \langle x | \mathcal {H} _ {\mathrm{A}} (| f | ^ {2}) x \rangle = \int_ {\mathsf {X} (\mathrm{A})} | f | ^ {2} d \mu_ {x}.
\]

由于 \(\mathcal{C}(\mathrm{X(A)})\) 在 \(\mathcal{L}^2 (\mathrm{X(A)},\mu_x)\) 中稠密，存在唯一的从 \(\mathrm{L}^2 (\mathrm{X(A)},\mu_x)\) 到 \(\mathrm{E}_x\) 中的等距线性映射，它延拓 \(\widetilde{\theta}_x\) ；我们把它记作 \(\theta_{x}\) 。映射 \(\theta_{x}\) 是满射的，因为它的像是闭的（I, p. 107 的引理 8）且包含 \(\mathrm{A}\cdot x\) 。因而它是一个等距同构。

设 \(a \in \mathrm{A}\) 。令 \(g = \mathcal{G}_{\mathrm{A}}(a) \in \mathcal{C}(\mathsf{X}(\mathrm{A}))\) 。空间 \(\mathrm{E}_x\) 在 \(a\) 下稳定。对 \(f \in \mathcal{C}(\mathsf{X}(\mathrm{A}))\) ，我们有

\[
\widetilde {\theta} _ {x} (m _ {g} (f)) = \mathcal {H} _ {\mathrm{A}} (\mathcal {G} _ {\mathrm{A}} (a) f) x = (a \circ \mathcal {H} _ {\mathrm{A}} (f)) x = a (\widetilde {\theta} _ {x} (f)),
\]

由此得出 \(\theta_x \circ m_g = a_x \circ \theta_x\) 。

推论. — 设 E 是复希尔伯特空间。设 A 是 \(\mathcal{L}(\mathrm{E})\) 的交换含幺星子代数，且具有循环向量 \(x\) 。设 \(\mu_x\) 是 \(x\) 关于 A 的谱测度。存在唯一的等距同构

\[
\theta_ {x} \colon \mathrm{L} ^ {2} (\mathsf {X} (\mathrm{A}), \mu_ {x}) \to \mathrm{E}
\]

使得 \(\theta_x(f) = \mathcal{H}_{\mathrm{A}}(f)\) 对每个函数 \(f \in \mathcal{C}(\mathsf{X}(\mathrm{A}))\) 成立。对 A 中每个 a，我们有 \(a \circ \theta_x = \theta_x \circ m_{\mathcal{G}_{\mathrm{A}}(a)}\) 。

事实上，这时我们有 \(E_{x}=E\) ，且 \(a_{x}=a\) （对所有 \(a\in A\) ）。

设 E 是复希尔伯特空间。设 \(x \in E\) ，A 是 \(\mathcal{L}(\mathrm{E})\) 的交换含幺星子代数。用 \(E_{x}\) 表示 E 的闭子空间 \(\overline{A \cdot x}\) 。对每个 \(a \in A\) ，我们有 \(a(\mathrm{E}_{x}) \subset \mathrm{E}_{x}\) 。用 \(A_{x}\) 表示 \(\mathcal{L}(\mathrm{E}_{x})\) 的交换含幺星子代数，它由 A 的元素经过渡到子空间而得的 \(E_{x}\) 的自同态组成。向量 x 是 \(A_{x} \subset \mathcal{L}(\mathrm{E}_{x})\) 的循环向量。

命题 9. --- 存在 E 的子集 C，使得 E 是诸空间 \(E_{x}\) （\(x \in C\)）的希尔伯特直和。若 E 是可数生成的，则集 C 是可数的。

设 \(\mathcal{O}\) 是由 E 的这样的子集 C 组成的集：诸子空间 \(\mathrm{E}_x\) （\(x\in \mathrm{C}\)）两两正交。集 \(\mathcal{O}\) 按包含关系有序，它是有限特征的（E, III, p. 34，定义 2），因为 C 属于 \(\mathcal{O}\) 当且仅当由 C 的两个元素组成的集属于 \(\mathcal{O}\) 。由 E, III, p. 35 的定理 1，存在 \(\mathcal{O}\) 的一个极大元 C。

设 F 是由诸子空间 \(E_{x}\) （\(x \in C\)）生成的 E 的闭子空间。只须证明 \(F^{\circ}\) 归结为 0 即可完成命题的证明。设 y 是 \(F^{\circ}\) 的元素。对每个 \(x \in C\) 及 A 中所有 a 与 b，我们有 \(\langle a(y) | b(x) \rangle = \langle y | a^{*}b(x) \rangle = 0\) ，因为 \(a^{*}b(x)\) 属于 \(E_{x}\) ，从而属于 F。由于元素 \(a(y)\) （相应地， \(b(x)\) ）生成 \(E_{y}\) （相应地， \(E_{x}\) ）的稠密子空间，我们有 \(E_{y} \subset E_{x}^{\circ}\) 。由于 C 在 O 中是极大的，这就意味着 \(E_{y}\) 归结为 0，从而 y = 0。

假设 E 是可数生成的。由于 \(E_{x}\) 对所有非零 \(x \in C\) 非零，任何使 E 成为诸空间 \(E_{x}\) （\(x \in C\)）的希尔伯特直和的集 C 都是可数的。

定理 1. --- 设 A 是 \(\mathcal{L}(\mathrm{E})\) 的交换含幺星子代数。存在局部紧拓扑空间 X、X 上的正测度 \(\mu\) 、等距同构 \(\theta\) ——从 \(\mathrm{L}^2 (\mathrm{X},\mu)\) 到 E 上——以及星代数等距态射 \(\pi\) ——从 A 到 \(\mathcal{C}_b(\mathrm{X})\) 中——使得对每个 \(a\in \mathrm{A}\) ，有

\[
a \circ \theta = \theta \circ m _ {\pi (a)}.
\]

由命题 9，存在 E 的子集 C，使得 E 是诸子空间 \(E_{x}\) （\(x \in C\)）的希尔伯特直和。设 X 为局部紧拓扑空间 \(\mathsf{X}(\mathsf{A}) \times \mathsf{C}\) ，其中 C 赋有离散拓扑。存在 X 上唯一的测度 \(\mu\) ，使得 \(\mu|(\mathsf{X}(\mathsf{A}) \times \{x\})\) 与 x 的谱测度 \(\mu_{x}\) 等同（对每个 \(x \in C\) ；INT, III, p. 65, § 2, n°1，命题 1）；这个测度是正的（参见同前）。于是有希尔伯特直和分解

\[
\mathrm{L} ^ {2} (\mathrm{X}, \mu) = \bigoplus_ {x \in \mathrm{C}} \mathrm{L} ^ {2} (\mathrm{X} (\mathrm{A}), \mu_ {x})
\]

（IV, p. 179 的命题 1, c) 应用于诸集 \(\mathsf{X}(\mathsf{A}) \times \{x\}\) （\(x \in \mathrm{C}\)））。因此存在唯一的等距 \(\theta: \mathrm{L}^2(\mathrm{X}, \mu) \to \mathrm{E}\) ，它与 \(\theta_x: \mathrm{L}^2(\mathsf{X}(\mathsf{A}), \mu_x) \to \mathrm{E}_x\) 在 \(\mathrm{L}^2(\mathsf{X}(\mathsf{A}), \mu_x)\) 上重合（对所有 \(x \in \mathrm{C}\) ）。

设 \(p \colon \mathrm{X} \to \mathrm{X}(\mathrm{A})\) 是典范投影；它是满射的，因而线性映射 \(p^* \colon \mathcal{C}_b(\mathrm{X}(\mathrm{A})) \to \mathcal{C}(\mathrm{X})\) ——由 \(f \mapsto f \circ p\) 所定义——是单射的。令 \(\pi = p^* \circ \mathcal{G}_{\mathrm{A}}\) ；这是从 A 到 \(\mathcal{C}_b(\mathrm{X})\) 中的星代数的单射态射；因而它是等距的（I, p. 112 的命题 9）。

设 \(a \in \mathrm{A}\) 。对每个 \(x \in \mathrm{C}\) ，我们有 \(a_x \circ \theta_x = \theta_x \circ m_{\mathcal{G}_A(a)}\) 。于是连续线性映射 \(a \circ \theta\) 与 \(\theta \circ m_{\pi(a)}\) 在诸子空间 \(\mathrm{L}^2(\mathrm{X}(\mathrm{A}), \mu_x)\) 上重合，因而相等。

推论（谱定理）. --- 设 E 是复希尔伯特空间。设 \(u \in \mathcal{L}(\mathrm{E})\) 是 E 的正规自同态。存在局部紧拓扑空间 X、X 上的正测度 \(\mu\) 、等距同构 \(\theta\) ——从 \(\mathrm{L}^{2}(\mathrm{X},\mu)\) 到 E 上——以及 X 上的连续有界函数 g，使得 \(u = \theta \circ m_{g} \circ \theta^{-1}\) 。

若 u 有循环向量 x，则可取 \(X = \mathrm{Sp}(u)\) ，并取 g 为 X 的恒等函数。

只须对 \(\mathcal{L}(\mathrm{E})\) 的、由 \(u\) 所生成的星子代数 A 应用定理 1（相应地，命题 8 的推论），并令 \(g = \pi(u)\) 。

这个命题把关于希尔伯特空间单个正规自同态的每个问题化归为关于乘法算子的类似问题，这往往使其研究得到简化（例如参见 IV, p. 325 的习题 19）。

\subsection{函数演算的连续性}

在本节中，我们考虑函数演算关于两个变量的连续性质。

对星代数 A 及 A 的正规元素 \(a\) ，以及对函数 \(f\) ——在 \(\mathbf{C}\) 的包含 \(\mathrm{Sp}_{\mathrm{A}}(a)\) 的子集上连续——我们用 \(f(a)\) 表示把 \(a\) 的连续函数演算应用于 \(f\) 在 \(a\) 的谱上的限制所得到的元素。

命题 10. --- 设 A 是含幺星代数。设 U 是 \(\mathbf{C}\) 中的相对紧开集。用 \(\Omega_{\mathrm{n}}\) 表示 A 中满足 \(\mathrm{Sp}_{\mathrm{A}}(a)\) 含于 U 的正规元素所成的集。映射 \((f,a)\mapsto f(a)\) ——从 \(\mathcal{C}(\overline{\mathrm{U}})\times \Omega_{\mathrm{n}}\) 到 A 中——是连续的。

用 \(q\) 表示映射 \((f, a) \mapsto f(a)\) ——从 \(\mathcal{C}(\overline{\mathrm{U}}) \times \Omega_{\mathrm{n}}\) 到 A 中。函数演算映射 \(f \mapsto f(a)\) （\(a \in \Omega_{\mathrm{n}}\)）所成的族在 \(\mathcal{L}(\mathcal{C}(\overline{\mathrm{U}}); \mathrm{A})\) 中是等度连续的，因为其中每个都是范数 \(\leqslant 1\) 的连续线性映射。因此，为证明断言，只须验证：对每个 \(f \in \mathcal{C}(\overline{\mathrm{U}})\) ，从 \(\Omega_{\mathrm{n}}\) 到 A 中的、由 \(a \mapsto f(a)\) 所定义的映射是连续的（TG, X, p. 13，推论 3）。

设 \(\mathcal{A}\) 是由 \(f\in\mathcal{C}(\overline{\mathrm{U}})\) （使映射 \(a\mapsto f(a)\) ——从 \(\Omega_{\mathrm{n}}\) 到 A 中——连续）组成的集；要证明的是 \(\mathcal{A}=\mathcal{C}(\overline{\mathrm{U}})\) 。

集 \(\mathcal{A}\) 是 \(\mathcal{C}(\overline{\mathrm{U}})\) 的含幺对合子代数。它包含 \(\overline{\mathrm{U}}\) 上的恒等函数，因而分离点。于是，它在 \(\mathcal{C}(\overline{\mathrm{U}})\) 中稠密（TG, X, p. 39，命题 7）。我们来证明它是闭的。

设 \((f_{\mathrm{n}})\) 是 \(\mathcal{A}\) 中收敛于 \(f \in \mathcal{C}(\overline{\mathrm{U}})\) 的序列。设 \(\varepsilon > 0\) ，并选取 \(n \in \mathbf{N}\) 使得 \(\|f - f_{n}\|_{\infty} \leqslant \varepsilon/4\) 。对每个 \((a_{1}, a_{2}) \in \Omega_{\mathrm{n}}^{2}\) ，我们有

\[
\begin{array}{c} \| f (a _ {1}) - f (a _ {2}) \| \leqslant 2 \| f - f _ {n} \| _ {\infty} + \| f _ {n} (a _ {1}) - f _ {n} (a _ {2}) \| \\ \leqslant \frac {\varepsilon}{2} + \| f _ {n} (a _ {1}) - f _ {n} (a _ {2}) \|. \end{array}
\]

由于 \(f_n\) 属于 \(\mathcal{A}\) ，存在 \(a_1\) 在 \(\Omega_n\) 中的邻域 V，使得 \(\| f_n(a_1) - f_n(a_2) \| \leqslant \varepsilon / 2\) 对所有 \(a_2 \in V\) 成立。于是，我们有 \(\| f(a_1) - f(a_2) \| \leqslant \varepsilon\) （对所有 \(a_2 \in V\) ）。因而映射 \(a \mapsto f(a)\) 在 \(a_1\) 处连续；断言得证。

推论 1. --- 设 A 是含幺星代数，\(A_{n}\) 是 A 的正规元素所成的集。给空间 \(\mathcal{C}(\mathbf{C})\) 赋以紧收敛拓扑。映射 \((f,a)\mapsto f(a)\) ——从 \(\mathcal{C}(\mathbf{C})\times\mathrm{A}_{\mathrm{n}}\) 到 A 中——是连续的。

设 \((f_{0}, a_{0}) \in \mathcal{C}(\mathbf{C}) \times \mathrm{A}_{\mathrm{n}}\) 。设 U 是 \(a_{0}\) 的谱的相对紧邻域。设 V 是 A 中满足 \(\mathrm{Sp}_{\mathrm{A}}(a) \subset \mathrm{U}\) 的正规元素 a 所成的集；它是 \(a_{0}\) 在 \(A_{n}\) 中的开邻域（I, p. 76，命题 10）。对每个 \(a \in V\) 及每个函数 \(f \in \mathcal{C}(\mathbf{C})\) ，我们有 \(f(a) = (f|\overline{\mathrm{U}})(a)\) 。由于映射 \(f \mapsto \|f|\overline{U}\|_{\infty}\) 是赋以紧收敛拓扑的空间 \(\mathcal{C}(\mathbf{C})\) 上的连续半范数，映射 \((f, a) \mapsto f(a)\) 在 \((f_{0}, a_{0})\) 处的连续性由命题 10 得出。

推论 2. --- 设 E 是复希尔伯特空间，\(\mathcal{L}(\mathrm{E})_{\mathrm{n}}\) 是 E 的正规自同态所成的集。给空间 \(\mathcal{C}(\mathbf{C})\) 赋以紧收敛拓扑。从 \(\mathcal{C}(\mathbf{C}) \times \mathcal{L}(\mathrm{E})_{\mathrm{n}}\) 到 \(\mathcal{L}(\mathrm{E})\) 中的、由 \((f, u) \mapsto f(u)\) 所定义的映射是连续的。

\section{广义函数与缓增广义函数}

在本节中，n 表示一个自然数。我们用 \(\mu\) 表示 \(\mathbf{R}^{n}\) 上的勒贝格测度，以及它在 \(\mathbf{R}^{n}\) 的每个局部紧子集上的限制。

我们用

\[
x \cdot y = \sum_ {i = 1} ^ {n} x _ {i} y _ {i}
\]

表示欧几里得空间 \(\mathbf{R}^{n}\) 上的数量积；欧几里得范数记作 \(x \mapsto \|x\|\) （TG, VI, p. 7）。我们回顾，群 \(\mathbf{R}^{n}\) 通过映射 \((x, y) \mapsto \exp(2i\pi x \cdot y)\) 与自身对偶，且勒贝格测度的对偶测度这时与勒贝格测度等同（II, p. 236 的推论 3）。我们用 F （相应地， \(\overline{F}\) ）表示 \(\mathbf{R}^{n}\) 上的傅里叶变换（相应地，傅里叶逆变换）（参见 II, p. 237 的 no. 9）。

对每个 \(\alpha = (\alpha_{i})_{1\leqslant i\leqslant n}\in \mathbf{N}^{n}\) ，我们用 \(X^{\alpha}\) 表示从 \(\mathbf{R}^n\) 到 \(\mathbf{R}\) 中的、由 \(x = (x_{i})_{1\leqslant i\leqslant n}\mapsto x^{\alpha} = \prod_{i = 1}^{n}x_{i}^{\alpha_{i}}\) 所定义的函数。

设 \(\alpha\) 与 \(\beta\) 是 \(\mathbf{N}^n\) 的元素。我们记

\[
| \alpha | = \sum_ {i = 1} ^ {n} \alpha_ {i}, \qquad \binom{\alpha}{\beta} = \prod_ {i = 1} ^ {n} \binom{\alpha_ {i}}{\beta_ {i}}.
\]

引理 1. --- 设 U 是 \(\mathbf{R}^{n}\) 的开子集。

\begin{enumerate}
\def\labelenumi{\alph{enumi})}
\item
  设 K 是 U 的紧子集，V 是 K 在 U 中的开邻域。存在 \(\varphi \in \mathcal{C}^{\infty}(\mathrm{U})\) ，其紧支撑含于 V，使得 \(0 \leqslant \varphi \leqslant 1\) ，且 \(\varphi(x) = 1\) 对所有 \(x \in \mathrm{K}\) 成立；
\item
  对 U 的每个局部有限开覆盖，存在一个由 \(C^{\infty}\) 类函数组成、从属于它的单位分解。
\end{enumerate}

这由 VAR, R1, p. 40, 5.3.6 得出。

命题 1. --- 设 U 是 \(\mathbf{R}^{n}\) 的开子集，\(k\in\mathbf{N}\) 。设 E、F、G 是拓扑向量空间，b: E×F→G 是连续双线性映射。设 f 与 g 分别是 U 到 E 与 F 中的 \(\mathbf{C}^{k}\) 类函数。那么映射 h: x↦b(f(x),g(x)) 在 U 上是 \(\mathbf{C}^{k}\) 类的；此外，对每个 \(\alpha\in\mathbf{N}^{n}\) （满足 \(|\alpha|\leqslant k\) ）以及

对每个 \(x \in \mathrm{U}\) ，我们有

\[
\partial^{\alpha}h(x) = \sum_{\substack{\beta \in \mathbf{N}^{n}\\ \beta \leqslant \alpha}}\binom {\alpha}{\beta}b\big(\partial^{\beta}f(x),\partial^{\alpha -\beta}g(x)\big).
\]

映射 h 是映射 \((f,g)\) : \(U \rightarrow E \times F\) （它是 \(C^{k}\) 类的）与映射 b: \(E \times F \rightarrow G\) （它是 \(C^{\infty}\) 类的）的复合。因而它是 \(C^{k}\) 类的。其偏导数的表达式由莱布尼茨公式得出（FVR, I, p. 28，命题 2，参见 A, III, p. 122，推论）。

我们回顾，分离的局部凸拓扑向量空间若为桶型空间且每个有界子集都相对紧，则称为蒙泰尔空间（EVT, IV, p. 18，定义 4）。

我们还回顾，若 E 与 F 是局部凸空间，且 E 是有界型的（EVT, III, p. 12，定义 1），则线性映射 \(u\colon\mathrm{E}\to\mathrm{F}\) 连续当且仅当 E 的每个有界子集在 \(u\) 下的像在 F 中有界（EVT, III, p. 11，命题 1, (iiibis)；当 F 是可赋范时，一般情形可形式地由此推出，参见 EVT, II, p. 7，命题 5, b)）。

\subsection{积分号下求导}

我们固定一个局部紧拓扑空间 X 和 X 上的一个测度 \(\nu\) 。设 E 是巴拿赫空间。

命题 2. --- 设 I 是 \(\mathbf{R}\) 的区间，f 是从 X × I 到 E 中的映射，满足

\begin{enumerate}
\def\labelenumi{(\roman{enumi})}
\item
  对每个 \(t \in I\) ，从 X 到 E 中的映射 \(x \mapsto f(x, t)\) 是 \(\nu\)-可积的；
\item
  对每个 \(x \in X\) ，从 I 到 E 中的映射 \(t \mapsto f(x, t)\) 有导数，记作 \(t \mapsto f'(x, t)\) ；
\item
  存在 X 上的正的 \(\nu\)-可积函数 g，使得 \(\|f'(x,t)\| \leqslant g(x)\) 对所有 \((x,t) \in \mathrm{X} \times \mathrm{I}\) 成立。
\end{enumerate}

那么由

\[
\mathrm{F} (t) = \int_ {\mathrm{X}} f (x, t) d \nu (x)
\]

定义的从 I 到 E 中的映射 F 在 I 中可微，且对每个 \(t \in \mathrm{I}\) ，我们有

\[
\mathrm{F} ^ {\prime} (t) = \int_ {\mathrm{X}} f ^ {\prime} (x, t) d \nu (x).
\]

设 \(t_0 \in \mathrm{I}\) ，\(\mathrm{J}\) 是 \(\mathbf{R}\) 的区间，含于 \(\mathrm{I}\) ，且是 \(t_0\) 在 \(\mathrm{I}\) 中的邻域。设 \(h: \mathrm{X} \times \mathrm{J} \to \mathrm{E}\) 是由 \((x, t) \mapsto (f(x, t) - f(x, t_0)) / (t - t_0)\) （对 \((x, t) \in \mathrm{X} \times (\mathrm{J} - \{t_0\})\) ）以及 \((x, t_0) \mapsto f'(x, t_0)\) （对 \(x \in \mathrm{X}\) ）所定义的映射。设 \(x \in \mathrm{X}\) 。我们有 \(\|h(x, t_0)\| \leqslant g(x)\) ，且对所有 \(t \neq t_0\) ，可得 \(\|h(x, t)\| \leqslant g(x)\) 。依据 FVR, I, p. 23 的定理 2。由导数的定义，命题由应用于映射 \(h\) 的 INT, IV, p. 144, § 4, no. 3 的推论 1 得出。

我们采用 VAR, R1, 2.4, p. 19 中关于偏导数的记号。

推论 1. --- 设 U 是 \(\mathbf{R}^{n}\) 的开子集。设 \(k \in \mathbf{N}\) ，\(f\) 是从 \(\mathrm{X} \times \mathrm{U}\) 到 E 中的映射，满足下列条件：

\begin{enumerate}
\def\labelenumi{(\roman{enumi})}
\item
  对每个 \(t \in U\) ，从 X 到 E 中的映射 \(x \mapsto f(x, t)\) 是 \(\nu\)-可积的；
\item
  对每个 \(x \in X\) ，从 U 到 E 中的映射 \(t \mapsto f(x, t)\) 是 \(C^{k}\) 类的，其偏导数（记作 \(\partial^{\alpha}f(x, t)\) ，对每个 \(\alpha \in \mathbf{N}^{n}\) 满足 \(|\alpha| \leqslant k\) ）；
\item
  存在 \(\nu\)-可积的函数 \(g\) （X 上的），使得对每个 \(\alpha \in \mathbf{N}^n\) （满足 \(|\alpha| \leqslant k\) ）以及每个 \((x,t) \in \mathrm{X} \times \mathrm{U}\) ，有 \(\|\partial^\alpha f(x,t)\| \leqslant g(x)\) 。
\end{enumerate}

那么由

\[
\mathrm{F} (t) = \int_ {\mathrm{X}} f (x, t) d \nu (x)
\]

定义的从 U 到 E 中的映射 F 在 U 中是 \(C^{k}\) 类的，且对每个 \(\alpha\in \mathbf{N}^{n}\) （满足 \(|\alpha|\leqslant k\) ）及每个 \(t\in U\) ，我们有

\[
\partial^ {\alpha} \mathrm{F} (t) = \int_ {\mathrm{X}} \partial^ {\alpha} f (x, t) d \nu (x).
\]

这由命题对 k 作归纳得出。

推论 2. --- 设 k 是自然数。设 \(f \in \mathcal{L}^1(\mathbf{R}^n, \mu)\) ，\(g \in \mathcal{C}^k(\mathbf{R}^n)\) 。假设对每个 \(\alpha \in \mathbf{N}^n\) （满足 \(|\alpha| \leqslant k\) ），函数 \(\partial^\alpha g\) 都有界。那么卷积 \(f * g\) 属于 \(\mathcal{C}^k(\mathbf{R}^n)\) ，且对每个 \(\alpha\) （满足 \(|\alpha| \leqslant k\) ），有 \(\partial^\alpha(f * g) = f * \partial^\alpha g\) 。

我们可以对赋勒贝格测度的空间 \(X = \mathbf{R}^{n}\) 以及由 \((x, t) \mapsto f(x)g(t - x)\) 定义的从

\(\mathbf{R}^{n} \times \mathbf{R}^{n}\) 到 \(\mathbf{C}\) 中的映射 h 应用推论 1；事实上，对每个 \(\alpha \in \mathbf{N}^{n}\) （满足 \(|\alpha| \leqslant k\) ），我们有不等式

\[
| \partial^ {\alpha} h (x, t) | \leqslant \Bigl (\sup _ {| \beta | \leqslant k} \sup _ {y \in \mathbf {R} ^ {n}} | \partial^ {\beta} g (y) | \Bigr) | f (x) |
\]

其右端是 \(R^{n}\) 上的可积函数。

\subsection{\texorpdfstring{\(\mathbf{R}^{n}\) 与 \(\mathbf{Z}^{n}\) 中的可积性判据}{\textbackslash mathbf\{R\}\^{}\{n\} 与 \textbackslash mathbf\{Z\}\^{}\{n\} 中的可积性判据}}

命题 3. --- 设 N 是 \(\mathbf{R}^{n}\) 上的一个范数，\(r\) 是实数，\(p \in [1, +\infty[\) 。

\begin{enumerate}
\def\labelenumi{\alph{enumi})}
\tightlist
\item
  函数 \((1 + \mathrm{N})^{r}\) 属于 \(\mathcal{L}^{p}(\mathbf{R}^{n}, \mu)\) 当且仅当 \(rp < -n\)；
\end{enumerate}

a') 在 \(\mathbf{Z}^{n}\) 上，函数 \((1+N)^{r}\) 的限制属于 \(\ell^{p}(\mathbf{Z}^{n})\) 当且仅当 \(rp < -n\)；

\begin{enumerate}
\def\labelenumi{\alph{enumi})}
\setcounter{enumi}{1}
\tightlist
\item
  设 V 是 \(\mathbf{R}^{n}\) 中 0 的有界可测邻域。函数 \(\mathbf{N}^{r}\) 属于 \(\mathcal{L}^{p}(\mathbf{R}^{n}-\mathbf{V},\mu)\) 当且仅当 \(rp<-n\) 。
\end{enumerate}

由 \(\mathbf{R}^{n}\) 上范数的等价性（EVT I, p. 14，定理 2 及 TG, IX, p. 32，命题 8），可以假设范数 N 由 \(\mathrm{N}(x)=\sup|x_{i}|\) 给出，其中 \(x=(x_{i})\in\mathbf{R}^{n}\) 。用 B 表示 \(\mathbf{R}^{n}\) 对此范数的单位球。

设 V 是 \(\mathbf{R}^n\) 中 0 的有界可测邻域。函数 \(N^r\) 在 \((\mathrm{B - V})\cup (\mathrm{V - B})\) 上连续且有界，这表明断言 \(b)\) 成立当且仅当 \(V = B\) 时它成立，以下我们就这样假设。函数 \((1 + N)^r\) 在 B 上连续且有界，且满足不等式 \(N^r\leqslant (1 + N)^r\leqslant 2^r N^r\) （在 \(\mathbf{R}^n -\mathbf{B}\) 上），由此可看出断言 \(a)\) 与 \(b)\) 等价。

我们来证明 V = B 时的 b)。可以假设 p = 1。\(r > 0\) 的情形是初等的，故设 \(r \leqslant 0\) 。对每个整数 \(j \geqslant 1\) ，可积集

\[
\mathrm{C} _ {j} = \{x \in \mathbf {R} ^ {n} \mid 2 ^ {j - 1} \leqslant \mathrm{N} (x) <   2 ^ {j} \}
\]

的测度为 \(2^{nj}(2^n - 1)\) 。诸集合 \((\mathrm{C}_j)_{j \geqslant 1}\) 构成 \(\mathbf{R}^n - \mathbf{B}\) 的一个划分。依据 INT, V, p. 4, § 1, no. 1，推论，我们因此有

\[
\begin{array}{c} \int_ {\mathbf {R} ^ {n} - \mathrm{B}} ^ {*} \mathrm{N} ^ {r} d \mu = \sum_ {j \geqslant 1} \int_ {\mathrm{C} _ {j}} ^ {*} \mathrm{N} ^ {r} d \mu \\ \leqslant \sum_ {j \geqslant 1} 2 ^ {n + n j} 2 ^ {r (j - 1)} = 2 ^ {n - r} \sum_ {j \geqslant 1} 2 ^ {(n + r) j} \end{array}
\]

它在 \(r < -n\) 时有限。另一方面，我们有（同前）

\[
\int_ {\mathbf {R} ^ {n} - \mathrm{B}} ^ {*} \mathrm{N} ^ {r} d \mu \geqslant \sum_ {j \geqslant 1} 2 ^ {r j} 2 ^ {n j} (2 ^ {n} - 1) = (2 ^ {n} - 1) \sum_ {j \geqslant 1} 2 ^ {(r + n) j}
\]

它在 \(r \geqslant -n\) 时无穷。

我们类似地证明 \(a'\) )，所考虑的集合 \(\mathrm{C}_{j} \cap \mathbf{Z}^{n}\) 覆盖 \(\mathbf{Z}^{n} - \{0\}\) 并满足

\[
\operatorname{Card} \left(\mathrm{C} _ {j} \cap \mathbf {Z} ^ {n}\right) = (2 ^ {j + 1} - 1) ^ {n} - (2 ^ {j} - 1) ^ {n},
\]

该数属于区间 \([(1-2^{-n})(2^{j+1}-1)^{n},2^{n(j+1)}]\) 。

\subsection{测试函数}

在本章中，U 表示 \(\mathbf{R}^{n}\) 的开子集。对 \(p \in [1, +\infty]\) ，我们将写 \(\mathcal{L}^{p}(U)\) 与 \(L^{p}(U)\) ，而不写 \(\mathcal{L}^{p}(U, \mu)\) 与 \(L^{p}(U, \mu)\) 。我们把属于 \(\mathcal{L}^{p}(U)\) 的连续函数与其在 \(L^{p}(U)\) 中的像等同起来。

给 U 上的复值无穷次可微函数空间 \(\mathcal{C}^{\infty}(\mathrm{U})\) 赋以由半范数族 \(p_{\alpha,K}\) ——对 \(\alpha \in \mathbf{N}^{n}\) 和 \(K \subset U\) 定义为

\[
p _ {\alpha , \mathrm{K}} (\varphi) = \sup _ {x \in \mathrm{K}} | \partial^ {\alpha} \varphi (x) |.
\]

该空间是完备的（参见 EVT, III, p. 9，例 b））。

对 U 的每个紧子集 K，我们用 \(\mathcal{C}_{\mathrm{K}}^{\infty}(\mathrm{U})\) 表示 \(\mathcal{C}^{\infty}(\mathrm{U})\) 中由支集含于 K 的函数构成的子空间。空间 \(\mathcal{C}_{\mathrm{K}}^{\infty}(\mathrm{U})\) 赋以由 \(\mathcal{C}^{\infty}(\mathrm{U})\) 的拓扑诱导的拓扑。它是一个弗雷歇空间，定义其拓扑的可数半范数族即族 \((p_{\alpha,\mathrm{K}})_{\alpha\in\mathbf{N}^{n}}\) 。特别地，它是有界型空间（EVT, III, p. 12，命题 2）。

我们用 \(\mathcal{D}(\mathrm{U})\) 表示向量空间 \(\mathcal{K}(\mathrm{U})\cap \mathcal{C}^{\infty}(\mathrm{U})\) ——由 U 中具有紧支集的 \(\mathrm{C}^{\infty}\) 类函数所构成。我们称 \(\mathcal{D}(\mathrm{U})\) 为 U 上的测试函数空间。空间 \(\mathcal{D}(\mathrm{U})\) 是诸空间 \(\mathcal{C}_{\mathrm{K}}^{\infty}(\mathrm{U})\) 的归纳极限，并赋以相应的局部凸归纳极限拓扑（EVT, II, p. 31）。

该空间在 EVT, III, p. 9 中记作 \(\mathcal{C}_{\circ}^{\infty}(U)\) 。

设 \((\mathrm{K}_{m})\) 是 U 的紧子集的递增序列，其内部构成 U 的覆盖。那么空间 \(\mathcal{D}(\mathrm{U})\) 是诸空间 \(\mathcal{C}_{\mathrm{K}_{m}}^{\infty}(\mathrm{U})\) 的严格归纳极限（EVT, III, p. 9）。因此它是完备空间（EVT, II, p. 35，命题 9）。\(\mathcal{D}(\mathrm{U})\) 的每个有界子集都含于某个子空间 \(\mathcal{C}_{\mathrm{K}_{m}}^{\infty}(\mathrm{U})\) 之中（EVT, III, p. 5，命题 6）。这就是说，\(\mathrm{B}\subset\mathcal{D}(\mathrm{U})\) 有界当且仅当存在 U 的紧子集 K 以及族 \((\mathrm{M}_{\alpha})_{\alpha\in\mathbf{N}^{n}}\) （\(R_{+}\) 中的），使得 B 含于诸函数 \(\varphi\in\mathcal{C}_{\mathrm{K}}^{\infty}(\mathrm{U})\) ——满足

\[
p _ {\alpha , \mathrm{K}} (\varphi) \leqslant \mathrm{M} _ {\alpha}
\]

——所成之集（对所有 \(\alpha\in\mathbf{N}^{n}\) ）。

空间 \(\mathcal{D}(\mathrm{U})\) 是有界型空间（EVT, III, p. 12，例 3）和蒙泰尔空间（EVT, IV, p. 18，例 4）。

设 V 是含于 U 的 \(\mathbf{R}^{n}\) 的开子集。若 \(K \subset V\) 是紧的，则函数到 V 上的限制定义了从 \(\mathcal{C}_{\mathrm{K}}^{\infty}(\mathrm{U})\) 到 \(\mathcal{C}_{\mathrm{K}}^{\infty}(\mathrm{V})\) 上的连续满线性映射。定义在 V 上的函数用零延拓，诱导一个从 \(\mathcal{D}(\mathrm{V})\) 到 \(\mathcal{D}(\mathrm{U})\) 中的单连续线性映射，称为典范映射。

注. --- 我们类似地定义 U 上的实值测试函数空间 \(\mathcal{D}_{\mathbf{R}}(\mathrm{U})\) 。使得 \(z \otimes \varphi \mapsto z\varphi\) （对所有 \(z \in \mathbf{C}\) 及每个 \(\varphi \in \mathcal{D}_{\mathbf{R}}(\mathrm{U})\) ）的线性映射是 \(\mathbf{C} \otimes \mathcal{D}_{\mathbf{R}}(\mathrm{U})\) 到 \(\mathcal{D}(\mathrm{U})\) 上的同构。

引理 2. --- 设 \(\mathcal{U}\) 是 U 的开覆盖。存在比 \(\mathcal{U}\) 更细的、由相对紧集合组成的 U 的局部有限开覆盖。

由于 U 是局部紧的，存在 \(\mathcal{V}\) ——U 的覆盖，比 \(\mathcal{U}\) 更细，由在 U 中相对紧的开集组成。由于 U 是仿紧的（TG, IX, p. 51，定理 4），存在局部有限开覆盖 \(\mathcal{W}\) ，比 \(\mathcal{V}\) 更细。于是 \(\mathcal{W}\) 比 \(\mathcal{U}\) 更细，且由相对紧的开集组成。

引理 3. --- 设 \(f \in \mathcal{K}(\mathbf{R}^n)\) ，V 是 f 的支集 K 的开邻域。存在整数 \(m_0 \geqslant 1\) ，使得对每个 \(x \in K\) ，集合 V 包含以 x 为中心、半径为 \(m_0^{-1}\) 的球。对每个整数 \(m > m_0\) ，设 \(V_m\) 是以 0 为中心、半径为 \(m^{-1}\) 的闭球（在 \(\mathbf{R}^{n}\) 中），\(\varphi_{m}\) 是支集含于 \(V_{m}\) 且积分为 1 的正测试函数。令 \(f_{m} = \varphi_{m} * f\) 。

\begin{enumerate}
\def\labelenumi{\alph{enumi})}
\item
  有 \(f_{m}\in \mathcal{D}(\mathbf{R}^{n})\) ，且 \(f_{m}\) 的支集含于 V；
\item
  设 \(p \in [1, +\infty]\) 。序列 \((f_m)_{m > m_0}\) 收敛于 \(f\) （在 \(L^p(\mathbf{R}^n)\) 中）。
\end{enumerate}

依据 TG, IX, p. 14，注，我们有 \(d(\mathrm{K}, \mathbf{R}^{n}-\mathrm{V}) > 0\) 。设 \(m_{0} \geqslant 1\) 是满足 \(m_{0}^{-1} < d(\mathrm{K}, \mathbf{R}^{n}-\mathrm{V})\) 的整数；它满足所需的条件。对所有 \(m > m_{0}\) ，函数 \(f_{m}\) 连续（INT, VIII, p. 166, § 4, n° 5，命题 11），且其支集含于 \(\mathrm{K} + \mathrm{V}_{m}\) （INT, VIII, p. 126, § 1, n° 4，命题 5），从而含于 V。依据 IV, p. 198 的推论 2，我们有 \(f_{m} \in \mathcal{D}(\mathbf{R}^{n})\) 。若 \(p \neq +\infty\) ，则序列 ( \(f_{m}\) ) 收敛于 \(f\) （在 \(\mathrm{L}^{p}(\mathbf{R}^{n})\) 中；INT, VIII, p. 172, § 4, no. 7，命题 20）。\(^{(1)}\)

设 \(p = +\infty\) 。设 \(\varepsilon > 0\) 。函数 \(f\) 在 \(\mathbf{R}^n\) 上一致连续；因此存在整数 \(m_1 > m_0\) ，使得对每个 \(m \geqslant m_1\) 以及所有 \(x \in \mathbf{R}^n\) 和 \(y \in \mathrm{V}_m\) ，有 \(|f(x) - f(x - y)| < \varepsilon\) 。

对每个 \(m \geqslant m_{1}\) ，函数 \(\varphi_{m}\) 在 \(\mathrm{V}_{m}\) 外为零、是正的且积分为 1。因此我们推出不等式

\[
| f (x) - f _ {m} (x) | = \left| \int_ {\mathbf {R} ^ {n}} (f (x) - f (x - y)) \varphi_ {m} (y) d \mu (y) \right| \leqslant \varepsilon
\]

（对所有 \(x \in \mathbf{R}^{n}\) ），由此得结论。

命题 4. --- a) \(\mathcal{D}(\mathrm{U})\) 到 \(\mathcal{K}(\mathrm{U})\) 中的典范单射是连续的，且 \(\mathcal{D}(\mathrm{U})\) 在 \(\mathcal{K}(\mathrm{U})\) 中稠密；

\begin{enumerate}
\def\labelenumi{\alph{enumi})}
\setcounter{enumi}{1}
\tightlist
\item
  对每个 \(p \in [1, +\infty[\) ，\(\mathcal{D}(\mathrm{U})\) 到 \(\mathcal{L}^p(\mathrm{U})\) 中的典范单射是连续的，且 \(\mathcal{D}(\mathrm{U})\) 在 \(\mathrm{L}^p(\mathrm{U})\) 中稠密。
\end{enumerate}

\(\mathcal{K}(\mathrm{U})\) 上的每个连续半范数在 \(\mathcal{D}(\mathrm{U})\) 上连续，因而 \(\mathcal{D}(\mathrm{U})\) 到 \(\mathcal{K}(\mathrm{U})\) 中的典范单射是连续的。

设 \(f \in \mathcal{K}(\mathrm{U})\) 。存在序列 \((f_{m})_{m \in \mathbf{N}}\) （在 \(\mathcal{D}(\mathrm{U})\) 中），它在 U 上一致收敛于 f，且诸 \(f_{m}\) 的支集都含于 f 的支集的一个固定的相对紧邻域（引理 3）。因此序列 \((f_{m})\) 收敛于 f （在 \(\mathcal{K}(\mathrm{U})\) 中）。这就完成了 a) 的证明。

设 \(p \in [1, +\infty[\) 。设 \(f \in \mathcal{D}(U)\) ，K 是 f 的支集。我们于是有不等式 \(\mathrm{N}_{p}(f) \leqslant \mu(\mathrm{K})^{1/p} \sup_{x \in \mathrm{K}} |f(x)|\) ，因而 \(\mathcal{D}(U)\) 到 \(\mathcal{L}^{p}(U)\) 中的典范单射是连续的。断言 b) 的后半部分由 a) 得出。

两个测试函数的乘积仍是测试函数，故莱布尼茨公式（FVR, I, p. 28，命题 2）表明空间 \(\mathcal{D}(\mathrm{U})\) 是一个拓扑代数。

更一般地，设 \(f \in \mathcal{C}^{\infty}(\mathrm{U})\) 。那么线性映射 \(\varphi \mapsto f\varphi\) ——从 \(\mathcal{D}(\mathrm{U})\) 到 \(\mathcal{D}(\mathrm{U})\) 中——是连续的。事实上，对每个紧子集 \(\mathbf{K}\) （\(\mathrm{U}\) 的）及每个 \(\alpha \in \mathbf{N}^{n}\) ，我们有

\[
p _ {\alpha , \mathrm{K}} (f \varphi) \leqslant \sum_ {0 \leqslant \beta \leqslant \alpha} \binom {\alpha} {\beta} p _ {\beta , \mathrm{K}} (f) p _ {\alpha - \beta , \mathrm{K}} (\varphi)
\]

（参见同前）。

\subsection{广义函数}

我们沿用上一节的记号；于是 U 表示 \(R^{n}\) 的开子集。

定义 1. --- \(\mathcal{D}(\mathrm{U})\) 的对偶空间，赋以有界收敛拓扑，称为 U 上的广义函数空间。记作 \(\mathcal{D}'(\mathrm{U})\) 。

于是，U 上的广义函数就是 \(\mathcal{D}(\mathrm{U})\) 上的连续线性形式。

若 \(f\) 是 \(\mathcal{D}(\mathrm{U})\) 上的线性形式（特别地，若 \(f\) 是广义函数），且 \(\varphi\) 属于 \(\mathcal{D}(\mathrm{U})\) ，我们用 \(\langle f, \varphi \rangle\) 表示 \(f\) 在 \(\varphi\) 处的取值。

由于 \(\mathcal{D}(\mathrm{U})\) 是有界型的，空间 \(\mathcal{D}'(\mathrm{U})\) 是完备的（TVS, III, p. 24，推论 1）。由于 \(\mathcal{D}(\mathrm{U})\) 是蒙泰尔空间，\(\mathcal{D}'(\mathrm{U})\) 亦然（EVT, IV, p. 19，命题 9）。特别地，空间 \(\mathcal{D}'(\mathrm{U})\) 是自反的（EVT, IV, p. 16，定理 2）。

引理 4. --- 设 f 是从 \(\mathcal{D}(\mathrm{U})\) 到 \(\mathbf{C}\) 中的线性映射。那么 \(f\) 是广义函数当且仅当：对每个紧子集 \(\mathrm{K}\) （\(\mathrm{U}\) 的）以及每个族 \((\mathrm{M}_{\alpha})_{\alpha \in \mathbf{N}^{n}}\) （\(\mathbf{R}_{+}\) 中的），线性形式 \(f\) 在由这样的函数 \(\varphi \in \mathcal{C}_{\mathrm{K}}^{\infty}(\mathrm{U})\) ——对每个 \(\alpha \in \mathbf{N}^{n}\) 有

\[
\sup _ {x \in \mathrm{K}} | \partial^ {\alpha} \varphi (x) | \leqslant \mathrm{M} _ {\alpha}.
\]

——所成之集上有界。事实上，空间 \(\mathcal{D}(\mathrm{U})\) 是有界型的，且 \(\mathcal{D}(\mathrm{U})\) 的每个有界子集都含于陈述中所描述的有界集之一。

引理 5. --- 设 F 是 \(\mathcal{D}'(\mathrm{U})\) 上具有可数基或包含一个单有界集合的滤子。那么 F 收敛于某个广义函数当且仅当 \(\langle f, \varphi \rangle\) 对 U 上的每个测试函数 \(\varphi\) 依 F 收敛。

特别地，广义函数序列 \((f_{m})_{m\in\mathbf{N}}\) 收敛于广义函数 f 当且仅当：对每个 \(\varphi\in\mathcal{D}(\mathrm{U})\) ，序列 \((\langle f_{m},\varphi\rangle)_{m\in\mathbf{N}}\) 收敛于 \(\langle f,\varphi\rangle\) 。

由于 \(\mathcal{D}'(\mathrm{U})\) 是蒙泰尔空间，故为桶型空间；这由巴拿赫–斯坦豪斯定理得出（EVT, III, p. 26，定理 1 的推论 3）。

设 V 是含于 U 的开集。从 \(\mathcal{D}(\mathrm{V})\) 到 \(\mathcal{D}(\mathrm{U})\) 中的典范线性映射的转置是从 \(\mathcal{D}'(\mathrm{U})\) 到 \(\mathcal{D}'(\mathrm{V})\) 中的连续线性映射，称为广义函数从 U 到 V 的限制，记作 \(r_{\mathrm{VU}}\) ，有时也记作 \(r_{\mathrm{V,U}}\) 。

有 \(r_{\mathrm{VV}} = 1_{\mathcal{D}'(\mathrm{V})}\) 。若 \(W \subset V \subset U\) 是 U 的开子集，则我们有 \(r_{\mathrm{WV}} \circ r_{\mathrm{VU}} = r_{\mathrm{WU}}\) 。换言之，若用 \(\mathcal{T}\) 表示 U 的拓扑，则投射系 \(\underline{\mathcal{D}'}(\mathrm{U}) = ((\mathcal{D}'(\mathrm{V}))_{\mathrm{V} \in \mathcal{T}}, (r_{\mathrm{WV}}))\) 是 U 上取值于局部凸拓扑向量空间的结构种的预层（TA, I, p. 42，定义 1 及 p. 66, §10）。

命题 5. --- 预层 \(\underline{\mathcal{D}}^{\prime}(\mathrm{U})\) 是一个层。

回顾（TA, I, p. 43，定义 2），这意味着对 U 的每个开子集 V 及 V 的每个开覆盖 \((\mathrm{V}_{i})_{i\in\mathrm{I}}\) ，下列条件满足：

\begin{enumerate}
\def\labelenumi{(\roman{enumi})}
\item
  映射 \((r_{\mathrm{V}_i\mathrm{V}})_{i\in \mathrm{I}}\colon \mathcal{D}'(\mathrm{V})\to \prod_{i\in \mathrm{I}}\mathcal{D}'(\mathrm{V}_i)\) 是单射；
\item
  对每个族 \((f_i)\in \prod_{i\in \mathrm{I}}\mathcal{D}'(\mathrm{V}_i)\) ——满足
\end{enumerate}

\[
r _ {(\mathrm{V} _ {i} \cap \mathrm{V} _ {i ^ {\prime}}) \mathrm{V} _ {i}} (f _ {i}) = r _ {(\mathrm{V} _ {i} \cap \mathrm{V} _ {i ^ {\prime}}) \mathrm{V} _ {i ^ {\prime}}} (f _ {i ^ {\prime}})
\]

——对每对 \((i, i') \in \mathrm{I} \times \mathrm{I}\) 都成立，则存在广义函数 \(f \in \mathcal{D}'(\mathrm{V})\) ，使得对每个 \(i \in I\) ，有 \(r_{\mathrm{V}_{i}\mathrm{V}}(f) = f_{i}\) 。

设 V 是 U 的开子集，\(\mathcal{V} = (\mathrm{V}_{i})_{i \in \mathrm{I}}\) 是 V 的开覆盖。设 \((\mathrm{W}_{j})_{j \in \mathrm{J}}\) 是 V 的开覆盖，它局部有限、比 \(\mathcal{V}\) 更细，且由相对紧的部分组成（IV, p. 201，引理 2）。取定单位分解 \((\varphi_{j})_{j \in \mathrm{J}}\) ——从属于覆盖 \((\mathrm{W}_{j})_{j \in \mathrm{J}}\) ——使得 \(\varphi_{j}\) 的支集含于 \(W_{j}\) （对所有 \(j \in J\) ；IV, p. 196，引理 1）。

设 \(\varphi\in\mathcal{D}(V)\) 。由于这样的 \(j\in J\) ——使 \(W_{j}\) 与 \(\varphi\) 的支集相交——组成的集合是有限的（TG, I, §9, no. 1, p. 59），我们有

\[
\varphi = \sum_ {j \in \mathrm{J}} \varphi \varphi_ {j}
\]

其中和式只有有限多个非零项。

我们来证明 (i)。设 \(f \in \mathcal{D}'(\mathrm{V})\) 满足 \(r_{\mathrm{V}_i\mathrm{V}}(f) = 0\) （对所有 \(i \in \mathrm{I}\) ）。这意味着 \(\langle f, \varphi \rangle = 0\) 对每个测试函数 \(\varphi\) （其支集含于某个开集 \(\mathrm{V}_i\) ）成立。若 \(\varphi\) 的支集含于某个开集 \(\mathrm{W}_j\) ，则更是如此。但这时，对每个 \(\varphi \in \mathcal{D}(\mathrm{V})\) ，我们有

\[
\langle f, \varphi \rangle = \langle f, \sum_ {j \in \mathrm{J}} \varphi \varphi_ {j} \rangle = \sum_ {j \in \mathrm{J}} \langle f, \varphi \varphi_ {j} \rangle = 0,
\]

因而 f = 0。

我们来证明 (ii)。设 \((f_{i})_{i\in\mathrm{I}}\) 是这样的族：\(f_{i}\in\mathcal{D}^{\prime}(\mathrm{V}_{i})\) （对所有 \(i\in I\) ），且 \(r_{\mathrm{V}_{i}\cap\mathrm{V}_{i^{\prime}},\mathrm{V}_{i}}(f_{i})=r_{\mathrm{V}_{i}\cap\mathrm{V}_{i^{\prime}},\mathrm{V}_{i^{\prime}}}(f_{i^{\prime}})\) 对 I 中所有 i 与 \(i^{\prime}\) 成立。设 \(\iota:J\to I\) 是满足 \(\mathrm{W}_{j}\subset\mathrm{V}_{\iota(j)}\) （对所有 \(j\in J\) ）的映射。对每个 \(j\in J\) ，令 \(\widetilde{f}_{j}=r_{\mathrm{W}_{j},\mathrm{V}_{\iota(j)}}(f_{\iota(j)})\) 。我们于是有

\[
\begin{array}{r l} & r _ {\mathrm{W} _ {j} \cap \mathrm{W} _ {j ^ {\prime}}, \mathrm{W} _ {j}} (\widetilde {f} _ {j}) = r _ {\mathrm{W} _ {j} \cap \mathrm{W} _ {j ^ {\prime}}, \mathrm{W} _ {j}} (r _ {\mathrm{W} _ {j}, \mathrm{V} _ {\iota (j)}} (f _ {\iota (j)})) \\ & \qquad = r _ {\mathrm{W} _ {j} \cap \mathrm{W} _ {j ^ {\prime}}, \mathrm{V} _ {\iota (j)}} (f _ {\iota (j)}) \\ & \qquad = r _ {\mathrm{W} _ {j} \cap \mathrm{W} _ {j ^ {\prime}}, \mathrm{V} _ {\iota (j)} \cap \mathrm{V} _ {\iota (j ^ {\prime})}} \circ r _ {\mathrm{V} _ {\iota (j)} \cap \mathrm{V} _ {\iota (j ^ {\prime})}, \mathrm{V} _ {\iota (j)}} (f _ {\iota (j)}). \end{array}
\]

交换 j 与 j' 的角色，并注意到

\[
r _ {\mathrm{V} _ {\iota (j)} \cap \mathrm{V} _ {\iota (j ^ {\prime})}, \mathrm{V} _ {\iota (j)}} (f _ {\iota (j)}) = r _ {\mathrm{V} _ {\iota (j)} \cap \mathrm{V} _ {\iota (j ^ {\prime})}, \mathrm{V} _ {\iota (j ^ {\prime})}} (f _ {\iota (j ^ {\prime})}),
\]

由假设成立，我们推出

\[
r _ {\mathrm{W} _ {j} \cap \mathrm{W} _ {j ^ {\prime}}, \mathrm{W} _ {j}} (\widetilde {f} _ {j}) = r _ {\mathrm{W} _ {j} \cap \mathrm{W} _ {j ^ {\prime}}, \mathrm{W} _ {j}} (\widetilde {f} _ {j ^ {\prime}})\tag{1}
\]

对所有 \((j, j') \in \mathrm{J}^{2}\) 成立。

对 \(\varphi \in \mathcal{D}(\mathrm{V})\) ，令

\[
\lambda (\varphi) = \sum_ {j \in \mathrm{J}} \langle \widetilde {f} _ {j}, \varphi \varphi_ {j} | \mathrm{W} _ {j} \rangle ,
\]

其中和式是有限的，因为只有有限多项可能非零。

映射 \(\lambda\) 是 \(\mathcal{D}(\mathrm{V})\) 上的线性形式。我们来证明它是广义函数。设 B 是 \(\mathcal{D}(\mathrm{V})\) 的有界子集，K 是 V 的紧子集且 \(B \subset \mathcal{C}_{K}^{\infty}(V)\) 。依据 TG, I, §9, no. 1, p. 59，存在 J 的有限子集 J' 使得

\[
\lambda (\varphi) = \sum_ {j \in J ^ {\prime}} \left\langle \widetilde {f} _ {j}, \varphi \varphi_ {j} \mid W _ {j} \right\rangle
\]

对所有 \(\varphi\in\mathrm{B}\) 成立。由于 \(\widetilde{f}_{j}\) 对每个 \(j\in\mathrm{J}^{\prime}\) 是广义函数，且 \(\mathcal{D}(\mathrm{V})\) 是拓扑代数，可知 \(\lambda\) 在 B 上有界，由此得结论（引理 4）。

设 \(j \in J\) ，\(\varphi \in \mathcal{D}(V)\) 的支集含于 \(W_{j}\) 。我们于是有

\[
\langle \lambda , \varphi \rangle = \sum_ {j ^ {\prime} \in \mathrm{J}} \left\langle \widetilde {f} _ {j ^ {\prime}}, \varphi \varphi_ {j ^ {\prime}} \right| \mathrm{W} _ {j} \rangle = \sum_ {j ^ {\prime} \in \mathrm{J}} \left\langle \widetilde {f} _ {j}, \varphi \varphi_ {j ^ {\prime}} \right| \mathrm{W} _ {j} \rangle
\]

这依据 (1)，因为 \(\varphi\varphi_{j'}\) 的支集含于 \(W_j \cap W_{j'}\) 。从而

\[
\langle \lambda , \varphi \rangle = \langle \widetilde {f} _ {j}, \sum_ {j ^ {\prime} \in \mathrm{J}} \varphi \varphi_ {j ^ {\prime}} | \mathrm{W} _ {j} \rangle = \langle \widetilde {f} _ {j}, \varphi | \mathrm{W} _ {j} \rangle ,
\]

因而 \(r_{\mathrm{W}_j\mathrm{V}}(\lambda) = \widetilde{f}_j\) 对所有 \(j\in \mathbf{J}\) 成立。

设 \(i \in I\) 。最后我们来证明 \(\lambda\) 在 \(V_i\) 上的限制与 \(f_i\) 一致。将条件 (i) 应用于 \(V_i\) 的由诸开集 \(W_j\) 组成的覆盖，只需验证对每个 \(j \in J\) ，\(\lambda\) 在 \(V_i \cap W_j\) 上的限制与 \(f_i\) 的限制一致。由前所述，只需验证 \(f_i\) 在 \(V_i \cap W_j\) 上的限制与 \(\widetilde{f}_j\) 的限制一致。而我们有

\[
\begin{array}{c} r _ {\mathrm{V} _ {i} \cap \mathrm{W} _ {j}, \mathrm{V} _ {i}} (f _ {i}) = r _ {\mathrm{V} _ {i} \cap \mathrm{W} _ {j}, \mathrm{V} _ {i} \cap \mathrm{V} _ {\iota (j)}} (r _ {\mathrm{V} _ {i} \cap \mathrm{V} _ {\iota (j)}, \mathrm{V} _ {i}} (f _ {i})) \\ = r _ {\mathrm{V} _ {i} \cap \mathrm{W} _ {j}, \mathrm{V} _ {i} \cap \mathrm{V} _ {\iota (j)}} (r _ {\mathrm{V} _ {i} \cap \mathrm{V} _ {\iota (j)}, \mathrm{V} _ {\iota (j)}} (f _ {\iota (j)})) = r _ {\mathrm{V} _ {i} \cap \mathrm{W} _ {j}, \mathrm{V} _ {\iota (j)}} (f _ {\iota (j)}), \end{array}
\]

其中我们用到了关于族 \((f_{i})\) 的假设。于是

\[
r _ {\mathrm{V} _ {i} \cap \mathrm{W} _ {j}, \mathrm{V} _ {\iota (j)}} (f _ {\iota (j)}) = r _ {\mathrm{V} _ {i} \cap \mathrm{W} _ {j}, \mathrm{W} _ {j}} (r _ {\mathrm{W} _ {j}, \mathrm{V} _ {\iota (j)}} (f _ {\iota (j)})) = r _ {\mathrm{V} _ {i} \cap \mathrm{W} _ {j}, \mathrm{W} _ {j}} (\widetilde {f} _ {j}),
\]

由此即可得出结论。

\subsection{函数解释为广义函数}

命题 6. --- 设 ν 是 U 上的测度。ν 在 D(U) 上的限制是一个广义函数，它为零当且仅当测度 ν 为零。

设 K 是 U 的紧子集。对每个函数 \(\varphi\in\mathcal{C}_{\mathrm{K}}^{\infty}(\mathrm{U})\) ，我们有 \(|\langle\nu,\varphi\rangle|\leqslant p_{0,\mathrm{K}}(\varphi)|\nu|(\mathrm{K})\) ，因而 \(\nu\) 在 \(\mathcal{D}(\mathrm{U})\) 上的限制是连续的。由于 \(\mathcal{D}(\mathrm{U})\) 在 \(\mathcal{K}(\mathrm{U})\) 中稠密，依据 IV, p. 202 的命题 4(a)，\(\nu\) 在 \(\mathcal{D}(\mathrm{U})\) 上的限制为零当且仅当 \(\nu\) 为零。

我们将把 U 上的复测度空间 \(\mathcal{M}(\mathrm{U};\mathbf{C})\) 与 \(\mathcal{D}'(\mathrm{U})\) 的一个子空间等同起来。我们还将把空间 \(\mathrm{L}_{\mathrm{loc}}^{1}(\mathrm{U})\) 与 \(\mathcal{D}'(\mathrm{U})\) 的一个子空间等同起来，所用的映射是由把 \(f\in \mathcal{L}_{\mathrm{loc}}^{1}(\mathrm{U})\) 映为测度 \(f\cdot \mu\) 的映射过渡到商而诱导的映射（INT, V, p. 44, § 5, no. 2，定义 2）。换言之，对 \(f\in \mathrm{L}_{\mathrm{loc}}^{1}(\mathrm{U})\) 与 \(\varphi \in \mathcal{D}(\mathrm{U})\) ，我们有

\[
\langle f, \varphi \rangle = \int_ {\mathrm{U}} f \varphi d \mu .
\]

特别地，对 \(p \in [1, +\infty]\) ，这使我们可以把空间 \(\mathrm{L}^p(\mathrm{U})\) 与 \(\mathcal{D}'(\mathrm{U})\) 的一个子空间等同起来（参见 INT, V, p. 43, § 5, no. 1）。

命题 7. --- 设 \(p \in [1, +\infty]\) 。\(\mathrm{L}^p(\mathrm{U})\) 到 \(\mathcal{D}'(\mathrm{U})\) 中的单射是连续的。

设 \(f \in \mathrm{L}^{p}(\mathrm{U})\) 。设 K 是 U 的紧子集，用 \(\varphi_{K}\) 表示它的特征函数。对每个支集含于 K 的测试函数 \(\varphi\) ，赫尔德不等式给出

\[
| \langle f, \varphi \rangle | = \left| \int_ {\mathrm{U}} f \varphi d \mu \right| \leqslant \mathrm{N} _ {p} (f) \mathrm{N} _ {q} (\varphi) \leqslant \mathrm{N} _ {p} (f) \mathrm{N} _ {q} (\varphi_ {\mathrm{K}}) p _ {0, \mathrm{K}} (\varphi)
\]

其中 q 是 p 的共轭指数。

特别地，可以把 \(\mathcal{D}(\mathrm{U})\) 与 \(\mathcal{D}'(\mathrm{U})\) 的一个子空间等同起来。

命题 8. --- 空间 \(\mathcal{D}(\mathrm{U})\) 在 \(\mathcal{D}'(\mathrm{U})\) 中稠密。

设 \(\lambda\) 是 \(\mathcal{D}'(\mathrm{U})\) 上在 \(\mathcal{D}(\mathrm{U})\) 上为零的线性形式。由于 \(\mathcal{D}(\mathrm{U})\) 是自反的，存在测试函数 \(\varphi \in \mathcal{D}(\mathrm{U})\) ，使得 \(\lambda(f) = \langle f, \varphi \rangle\) 对所有 \(f \in \mathcal{D}'(\mathrm{U})\) 成立。我们得到

\[
0 = \lambda (\overline {{\varphi}}) = \langle \overline {{\varphi}}, \varphi \rangle = \int_ {\mathrm{U}} | \varphi | ^ {2} d \mu ,
\]

因而 \(\varphi = 0\) 。命题由哈恩--巴拿赫定理得出（EVT, II, p. 49，推论 3 (ii)）。

对 \(h \in \mathcal{C}^{\infty}(\mathrm{U})\) ，连续线性映射 \(\varphi \mapsto h\varphi\) （\(\mathcal{D}(\mathrm{U})\) 到自身的）的转置是 \(\mathcal{D}'(\mathrm{U})\) 到自身的连续线性映射，记作 \(f \mapsto hf\) 。这个定义的合理性在于：若 \(f\) 是与测度 \(\nu\) （\(\mathrm{U}\) 上的）相伴的广义函数，则 \(hf\) 与测度 \(h \cdot \nu\) 相伴。事实上，对每个测试函数 \(\varphi \in \mathcal{D}(\mathrm{U})\) ，可以算出

\[
\langle h f, \varphi \rangle = \langle f, h \varphi \rangle = \int_ {\mathrm{U}} h \varphi d \nu = \int_ {\mathrm{U}} \varphi d (h \cdot \nu).
\]

\subsection{广义函数的求导}

设 \(\alpha \in \mathbf{N}^n\) 。线性映射 \(\varphi \mapsto \partial^\alpha \varphi\) 从 \(\mathcal{D}(\mathrm{U})\) 到 \(\mathcal{D}(\mathrm{U})\) 中是连续的。它的转置 \(^t\partial^\alpha\) 是从 \(\mathcal{D}'(\mathrm{U})\) 到 \(\mathcal{D}'(\mathrm{U})\) 中的连续线性映射（EVT, IV, p. 6，推论，\(b\) ））。

我们用 \(\partial^{\alpha}\) 表示连续线性映射 \((-1)^{|\alpha|}{}^{t}\partial^{\alpha}\) ——从 \(\mathcal{D}'(\mathrm{U})\) 到自身。

定义 2. --- 若 \(f \in \mathcal{D}'(\mathrm{U})\) 是广义函数，则称 \(\partial^{\alpha} f\) 为 \(\alpha\) 阶迭代偏导数（\(f\) 的）。

于是，按定义

\[
\langle \partial^ {\alpha} f, \varphi \rangle = (- 1) ^ {| \alpha |} \langle f, \partial^ {\alpha} \varphi \rangle
\]

（对每个函数 \(\varphi\in\mathcal{D}(\mathrm{U})\) ）。等式 \(\partial^{\alpha+\beta}=\partial^{\alpha}\circ\partial^{\beta}\) 对所有 \(\alpha\) 与 \(\beta\) （\(\mathbf{N}^{n}\) 中的）成立。

若 \(n=1\) ，也用 \(f'\) 记广义函数 \(f\in\mathcal{D}'(U)\) 的导数。该定义的合理性由下述引理给出。

引理 6. --- 设 k 是自然数。设 \(f \in \mathcal{C}^k(\mathrm{U})\) ，\(\lambda\) 是与 \(f\) 相伴的广义函数。对每个 \(\beta \in \mathbf{N}^n\) （满足 \(|\beta| \leqslant k\) ），广义函数 \(\partial^\beta\lambda\) 是与函数 \(\partial^\beta f\) 相伴的广义函数。

对 \(k\) 作归纳，只需在 \(\beta\) 满足 \(|\beta| = 1\) 时证明此性质，甚至可以假设 \(\beta = (0, \ldots, 0, 1)\) 。由于广义函数构成一个层（IV, p. 204，命题 5），只需在存在开集 \(V \subset \mathbf{R}^{n-1}\) 与开区间 \(I \subset \mathbf{R}\) 使得 \(U = V \times I\) 时验证断言。

对每个测试函数 \(\varphi\in\mathcal{D}(\mathrm{U})\) ，按定义有

\[
\langle \partial^ {\beta} \lambda , \varphi \rangle = - \int_ {\mathrm{U}} f (x) \partial^ {\beta} \varphi (x) d x = - \int_ {\mathrm{V}} \Bigl (\int_ {\mathrm{I}} f (y, t) \partial^ {\beta} \varphi (y, t) d t \Bigr) d y
\]

这里用到勒贝格--富比尼定理（INT, V, p. 96, § 8, n° 4，定理 1）。由分部积分（FVR, II, p. 10），有

\[
- \int_ {\mathrm{I}} f (y, t) \partial^ {\beta} \varphi (y, t) d t = \int_ {\mathrm{I}} \partial^ {\beta} f (y, t) \varphi (y, t) d t,
\]

因为 \(t \mapsto \varphi(y, t)\) 在 I 中有紧支集。于是得到

\[
\langle \partial^ {\beta} \lambda , \varphi \rangle = \int_ {V} \left(\int_ {I} \partial^ {\beta} f (y, t) \varphi (y, t) d t\right) d y = \langle \partial^ {\beta} f, \varphi \rangle .
\]

命题 9（莱布尼茨公式）. --- 设 f 是 U 上的广义函数，g 是 U 上的无穷次可微函数。设 \(\alpha \in \mathbf{N}^n\) 。我们有关系式

\[
\partial^ {\alpha} (f g) = \sum_ {\beta \leqslant \alpha} \binom {\alpha} {\beta} \partial^ {\beta} f   \partial^ {\alpha - \beta} g.
\]

如同 FVR, I, p. 28，命题 2 的证明那样对 \(|\alpha|\) 作归纳，只需考虑 \(|\alpha|=1\) 的情形。结论由计算

\[
\begin{array}{c} \langle \partial^ {\alpha} (f g), \varphi \rangle = \langle f g, - \partial^ {\alpha} \varphi \rangle = \langle f, - g \partial^ {\alpha} \varphi \rangle \\ = \langle f, - \partial^ {\alpha} (g \varphi) + \varphi \partial^ {\alpha} g \rangle = \langle g \partial^ {\alpha} f + f \partial^ {\alpha} g, \varphi \rangle \end{array}
\]

（对 \(\varphi\in\mathcal{D}(\mathrm{U})\) 成立）得出。

\subsection{施瓦茨函数}

我们用 \(\mathcal{S}(\mathbf{R}^{n})\) 表示无穷次可微函数 \(\varphi\) （在 \(\mathbf{R}^{n}\) 上、取复值）所成的空间，其中对所有 \(\alpha\) 与 \(\beta\) （属于 \(\mathbf{N}^{n}\) ），函数 \(\mathrm{X}^{\beta}\partial^{\alpha}\varphi\) 在 \(\mathbf{R}^{n}\) 上有界。我们给 \(\mathcal{S}(\mathbf{R}^{n})\) 赋以由可数半范数族 \((q_{\alpha,\beta})_{(\alpha,\beta)\in\mathbf{N}^{n}\times\mathbf{N}^{n}}\) 定义的局部凸拓扑，其中 \(q_{\alpha,\beta}\) 定义为

\[
q _ {\alpha , \beta} (\varphi) = \sup _ {x \in \mathbf {R} ^ {n}} | x ^ {\beta} \partial^ {\alpha} \varphi (x) | = \| \mathrm{X} ^ {\beta} \partial^ {\alpha} \varphi \| _ {\infty}
\]

（对 \(\varphi\in\mathcal{S}(\mathbf{R}^{n})\) ）。该拓扑是分离的。它也由半范数 \(\widetilde{q}_{\alpha,k}\) 定义，其中

\[
\widetilde {q} _ {\alpha , k} (\varphi) = \sup _ {x \in \mathbf {R} ^ {n}} \| x \| ^ {k} | (\partial^ {\alpha} \varphi) (x) |.
\]

对所有 \(\varphi\in\mathcal{S}(\mathbf{R}^{n})\) ，其中 \(k\in\mathbf{N}\) ，\(\alpha\in\mathbf{N}^{n}\) 。

我们称 \(\mathcal{S}(\mathbf{R}^{n})\) 为 \(\mathbf{R}^{n}\) 上的施瓦茨空间或施瓦茨函数空间。

注. --- 设 \(\varphi\in\mathcal{S}(\mathbf{R}^{n})\) 。对每个 \(k\in\mathbf{N}\) ，我们有

\[
\lim _ {\| x \| \to + \infty} \| x \| ^ {k} \varphi (x) = 0,
\]

因为函数 \(x \mapsto \|x\|^{k+1} \varphi(x)\) 有界。

例. --- 函数 \(\gamma_{n}\) ——在 \(\mathbf{R}^{n}\) 上由 \(\gamma_{n}(x)=\exp(-\|x\|^{2})\) 定义——属于 \(\mathcal{S}(\mathbf{R}^{n})\) 。事实上，对 k 作归纳可以证明：对每个整数 \(k\in \mathbf{N}\) ，存在多项式 \(P_{k}\in\mathbf{R}[X]\) 使得 \(\partial_{k}\gamma_{1}=P_{k}\gamma_{1}\) 。

于是，对所有 \(\alpha = (\alpha_{i})\in \mathbf{N}^{n}\) 与 \(\beta = (\beta_{i})\in \mathbf{N}^{n}\) ，以及每个 \(x = (x_{i})\in \mathbf{R}^{n}\) ，可得

\[
| (\mathrm{X} ^ {\beta} \partial^ {\alpha} \gamma_ {n}) (x) | = \prod_ {i = 1} ^ {n} | x _ {i} | ^ {\beta_ {i}} | \mathrm{P} _ {\alpha_ {i}} (x _ {i}) | \gamma_ {1} (x _ {i}),
\]

它在 x 于 \(\mathbf{R}^{n}\) 中变化时是有界量。

设 \(\alpha \in \mathbf{N}^n\) 与 \(\beta \in \mathbf{N}^n\) 。若 \(\varphi \in \mathcal{S}(\mathbf{R}^n)\) ，则 \(\mathrm{X}^{\beta}\partial^{\alpha}(\varphi)\) 仍是施瓦茨函数；如此定义的映射 \(\varphi \mapsto \mathrm{X}^{\beta}\partial^{\alpha}\varphi\) ——从 \(\mathcal{S}(\mathbf{R}^n)\) 到自身——是连续的。

空间 \(\mathcal{S}(\mathbf{R}^{n})\) 是一个拓扑代数。更确切地说，若 \(\varphi_{1}\) 与 \(\varphi_{2}\) 属于 \(\mathcal{S}(\mathbf{R}^{n})\) ，则 \(\varphi_{1}\varphi_{2}\) 是满足

\[
q _ {\alpha , \beta} (\varphi_ {1} \varphi_ {2}) \leqslant \sum_ {0 \leqslant \gamma \leqslant \alpha} \binom{\alpha}{\gamma} q _ {\gamma , \beta} (\varphi_ {1}) q _ {\alpha - \gamma , 0} (\varphi_ {2})\tag{2}
\]

（对所有 \(\alpha\) 与 \(\beta \in \mathbf{N}^n\) ；IV, p. 196，命题 1）的施瓦茨函数。

\(\mathcal{S}(\mathbf{R}^{n})\) 到空间 \(\mathcal{C}^{\infty}(\mathbf{R}^{n})\) （赋以 IV, p. 200 的 §3 中所述的拓扑）中的典范包含是连续的，因为

\[
\sup _ {x \in K} | \partial^ {\alpha} \varphi (x) | \leqslant q _ {\alpha , 0} (\varphi)
\]

对 \(\mathbf{R}^{n}\) 的每个紧子集 K、所有 \(\alpha \in \mathbf{N}^{n}\) 及每个施瓦茨函数 \(\varphi\) 都成立。

引理 7. --- 设 \(k \in \mathbf{N}\) ，\(\alpha \in \mathbf{N}^n\) 。对每个函数 \(\varphi \in \mathcal{S}(\mathbf{R}^n)\) 及每个实数 \(T > 0\)，我们有

\[
\widetilde {q} _ {\alpha , k} (\varphi) \leqslant \mathrm{T} ^ {k} \sup _ {\| x \| \leqslant \mathrm{T}} | \partial^ {\alpha} \varphi (x) | + \frac {1}{\mathrm{T}} \widetilde {q} _ {\alpha , k + 1} (\varphi).\tag{3}
\]

事实上，有

\[
\begin{array}{c} \widetilde {q} _ {\alpha , k} (\varphi) \leqslant \sup _ {\| x \| \leqslant T} \| x \| ^ {k} | \partial^ {\alpha} \varphi (x) | + \sup _ {\| x \| > T} \| x \| ^ {k} | \partial^ {\alpha} \varphi (x) | \\ \leqslant T ^ {k} \sup _ {\| x \| \leqslant T} | \partial^ {\alpha} \varphi (x) | + \frac {1}{T} \sup _ {\| x \| > T} \| x \| ^ {k + 1} | \partial^ {\alpha} \varphi (x) |. \end{array}
\]

命题 10. --- 设 B 是 \(\mathcal{S}(\mathbf{R}^{n})\) 的有界子集。\(\mathcal{S}(\mathbf{R}^{n})\) 在 B 上诱导的拓扑与 \(\mathcal{C}^{\infty}(\mathbf{R}^{n})\) 诱导的拓扑一致。

由于 \(\mathcal{S}(\mathbf{R}^{n})\) 到 \(\mathcal{C}^{\infty}(\mathbf{R}^{n})\) 中的包含是连续的，只需证明：对 \(\mathcal{S}(\mathbf{R}^{n})\) 的每个开子集 V，交集 \(V \cap B\) 对 \(\mathcal{C}^{\infty}(\mathbf{R}^{n})\) 诱导的拓扑而言在 B 中是开的。

设 V 是 \(\mathcal{S}(\mathbf{R}^{n})\) 的开子集。设 \(\varphi_{0} \in \mathrm{V} \cap \mathrm{B}\) 。存在有限集 I、族 \((\alpha_{i}, k_{i})_{i \in \mathrm{I}} \in (\mathbf{N}^{n} \times \mathbf{N})^{\mathrm{I}}\) 及实数 \(\varepsilon > 0\) ，使得 V 包含由 \(\varphi \in \mathcal{S}(\mathbf{R}^{n})\) ——满足

\[
\sup _ {i \in I} \widetilde {q} _ {\alpha_ {i}, k _ {i}} (\varphi - \varphi_ {0}) \leqslant \varepsilon .
\]

——所成之集。由于 B 在 \(\mathcal{S}(\mathbf{R}^{n})\) 中有界，存在 \(M > 0\) ，使得半范数 \(\widetilde{q}_{\alpha_{i},k_{i} + 1}\) （对 \(i\in I\) ）以 \(M\) 为界（在 \(B\) 上）。设 \(\delta >0\) 与 \(T > 0\) 为实数。依据不等式 (3)，只要 \(\varphi \in B\) 满足界

\[
\sup _ {i \in I} \sup _ {\| x \| \leqslant T} | \partial^ {\alpha_ {i}} (\varphi - \varphi_ {0}) | \leqslant \delta ,\tag{4}
\]

我们就有

\[
\sup _ {i \in I} \widetilde {q} _ {\alpha_ {i}, k _ {i}} (\varphi - \varphi_ {0}) \leqslant \delta T ^ {k} + \frac {2 M}{T}.
\]

取 \(T = \frac{4M}{\varepsilon}\) ，再取 \(\delta = \frac{\varepsilon}{2T^{k}}\) 。可以看出，\(V \cap B\) 包含 \(\varphi_{0}\) 在 B 中的、对 \(\mathcal{C}^{\infty}(\mathbf{R}^{n})\) 诱导拓扑而言由 (4) 定义的邻域。证明至此完成。

推论. --- 设 \((\varphi_{m})_{m\in\mathbf{N}}\) 是 \(\mathcal{S}(\mathbf{R}^{n})\) 中的有界序列。对每个函数 \(\varphi\in\mathcal{S}(\mathbf{R}^{n})\) ，下列断言等价：

\begin{enumerate}
\def\labelenumi{\alph{enumi})}
\item
  序列 \((\varphi_{m})\) 收敛于 \(\varphi\) （在 \(\mathcal{S}(\mathbf{R}^n)\) 中）；
\item
  序列 \((\varphi_{m})\) 收敛于 \(\varphi\) （在 \(\mathcal{C}^{\infty}(\mathbf{R}^{n})\) 中）。
\end{enumerate}

注. --- 序列 \((\varphi_{m})\) （在 \(\mathcal{C}^{\infty}(\mathbf{R}^{n})\) 中）收敛当且仅当：对每个 \(\alpha\in \mathbf{N}^{n}\) ，序列 \((\partial^{\alpha}\varphi_{m})\) 收敛于某个函数 \(\varphi^{(\alpha)}\) （在赋以紧收敛拓扑的 \(\mathcal{C}(\mathbf{R}^{n})\) 中）。这时我们有 \(\varphi^{(\alpha)}=\partial^{\alpha}\varphi\) ，且 \((\varphi_{m})\) 收敛于 \(\varphi^{(0)}\) 。

事实上，该条件是必要的。反之，若诸序列 \((\partial^{\alpha}\varphi_{m})\) 收敛于函数 \(\varphi^{(\alpha)}\) ——对所有 \(\alpha\in \mathbf{N}^{n}\) ——，则由 FVR, II, p. 2，定理 1 可得 \(\varphi^{(\alpha)}=\partial^{\alpha}\varphi^{(0)}\) ，这意味着序列 \((\varphi_{m})\) 收敛于 \(\varphi^{(0)}\) （在 \(\mathcal{C}^{\infty}(\mathbf{R}^{n})\) 中）。

命题 11. --- 空间 \(\mathcal{S}(\mathbf{R}^{n})\) 是弗雷歇空间和蒙泰尔空间。

由于空间 \(\mathcal{C}^{\infty}(\mathbf{R}^{n})\) 是完备的（EVT, III, p. 9，例 \(b\) ），命题 10 的推论表明 \(\mathcal{S}(\mathbf{R}^{n})\) 中的每个柯西序列都在 \(\mathcal{S}(\mathbf{R}^{n})\) 中收敛，因为它有界且在 \(\mathcal{C}^{\infty}(\mathbf{R}^{n})\) 中收敛。

于是空间 \(\mathcal{S}(\mathbf{R}^n)\) 是弗雷歇空间；特别地，它是桶型的（EVT, III, p. 25，命题 2 的推论）。设 B 是 \(\mathcal{S}(\mathbf{R}^n)\) 的有界子集，\((\varphi_m)_{m \in \mathbf{N}}\) 是取值于 B 的序列。由于 \(\mathcal{C}^\infty(\mathbf{R}^n)\) 是蒙泰尔空间（EVT, IV, p. 18，例 (4)），存在 \((\varphi_m)_{m \in \mathbf{N}}\) 的子序列在 \(\mathcal{C}^\infty(\mathbf{R}^n)\) 中收敛，从而在 \(\mathcal{S}(\mathbf{R}^n)\) 中收敛（命题 10）。因而 B 在 \(\mathcal{S}(\mathbf{R}^n)\) 中相对紧。由此得出 \(\mathcal{S}(\mathbf{R}^n)\) 是蒙泰尔空间。

\subsection{函数空间在施瓦茨函数空间中的包含}

命题 12. — 空间 \(\mathcal{D}(\mathbf{R}^{n})\) 含于 \(\mathcal{S}(\mathbf{R}^{n})\) ，且 \(\mathcal{D}(\mathbf{R}^{n})\) 到 \(\mathcal{S}(\mathbf{R}^{n})\) 中的包含是连续的且像稠密。

设 \(B \subset \mathcal{D}(\mathbf{R}^{n})\) 是有界子集，K 是 \(\mathbf{R}^{n}\) 的紧子集且 \(\mathrm{B} \subset \mathcal{C}_{\mathrm{K}}^{\infty}(\mathbf{R}^{n})\) 。设 \(\alpha \in \mathbf{N}^{n}\) 与 \(k \in \mathbf{N}\) 。对每个函数 \(\varphi \in B\) ，可得

\[
\widetilde {q} _ {\alpha , k} (\varphi) \leqslant \Bigl (\sup _ {x \in \mathrm{K}} \| x \| ^ {k} \Bigr) p _ {\alpha , \mathrm{K}} (\varphi),
\]

因而 B 在 \(\mathcal{S}(\mathbf{R}^{n})\) 中有界。包含的连续性则由空间 \(\mathcal{S}(\mathbf{R}^{n})\) 与 \(\mathcal{D}(\mathbf{R}^{n})\) 都是有界型的这一事实得出。

我们来证明 \(\mathcal{D}(\mathbf{R}^n)\) 在 \(\mathcal{S}(\mathbf{R}^n)\) 中稠密。设 B 是 \(\mathbf{R}^n\) 的单位球，\(\eta \in \mathcal{D}(\mathbf{R}^n)\) 是支集含于 2B 且满足 \(0 \leqslant \eta \leqslant 1\) 及 \(\eta(x) = 1\) （对所有 \(x \in \mathrm{B}\) ）的测试函数（IV, p. 196，引理 1, \(a\) ）。

设 \(\varphi \in \mathcal{S}(\mathbf{R}^n)\) 。对每个整数 \(m \geqslant 1\) 及每个 \(x \in \mathbf{R}^n\) ，令 \(\eta_m(x) = \eta(x/m)\) 。最后定义 \(\varphi_m = \eta_m\varphi\) ；我们有 \(\varphi_m \in \mathcal{D}(\mathbf{R}^n)\) 。

由于 \(\partial^{\alpha}\eta_{m} = m^{-|\alpha|}(\partial^{\alpha}\eta)(x / m)\) （对所有 \(\alpha \in \mathbf{N}^{n}\) 与 \(x\in \mathbf{R}^n\) ），由 IV, p. 210 的公式 (2) 推出序列 \((\varphi_{m})\) 在 \(\mathcal{S}(\mathbf{R}^n)\) 中有界。

设 C 是 \(R^{n}\) 的紧子集。序列 \((\varphi_{m})_{m\geqslant1}\) 收敛于 \(\varphi\) （在 \(\mathcal{C}_{\mathrm{K}}^{\infty}(\mathbf{R}^{n})\) 中），因为对所有充分大的 m，\(\varphi_{m}\) 在 C 上与 \(\varphi\) 一致。于是序列 \((\varphi_{m})\) 收敛于 \(\varphi\) （在 \(\mathcal{C}^{\infty}(\mathbf{R}^{n})\) 中），而 IV, p. 211 的命题 10 的推论使我们得出序列 \((\varphi_{m})\) 收敛于 \(\varphi\) （在 \(\mathcal{S}(\mathbf{R}^{n})\) 中）。

引理 8. --- 设 B 是 \(\mathcal{D}(\mathbf{R}^{n})\) 的有界子集。\(\mathcal{S}(\mathbf{R}^{n})\) 的拓扑在 B 上诱导的拓扑与 \(\mathcal{D}(\mathbf{R}^{n})\) 诱导的拓扑一致。

由于 \(\mathcal{D}(\mathbf{R}^{n})\) 到 \(\mathcal{S}(\mathbf{R}^{n})\) 中的包含是连续的，\(\mathcal{D}(\mathbf{R}^{n})\) 在 B 上诱导的拓扑比 \(\mathcal{S}(\mathbf{R}^{n})\) 诱导的拓扑更细。此外，存在 \(\mathbf{R}^{n}\) 的紧子集 K 使得 \(\mathrm{B} \subset \mathcal{C}_{\mathrm{K}}^{\infty}(\mathbf{R}^{n})\) 。于是对每个 \(\alpha \in \mathbf{N}^{n}\) ，有 \(p_{\alpha,K}(\varphi) \leqslant \widetilde{q}_{\alpha,0}(\varphi)\) ，这蕴含 \(\mathcal{S}(\mathbf{R}^{n})\) 诱导的拓扑比 \(\mathcal{D}(\mathbf{R}^{n})\) 诱导的拓扑更细。

命题 13. --- 设 \(p \in [1, +\infty]\) 。空间 \(\mathcal{S}(\mathbf{R}^{n})\) 含于 \(\mathcal{L}^{p}(\mathbf{R}^{n})\) ，且 \(\mathcal{S}(\mathbf{R}^{n})\) 到 \(\mathcal{L}^{p}(\mathbf{R}^{n})\) 中的典范单射是连续的。则 \(\mathcal{S}(\mathbf{R}^{n})\) 在 \(\mathrm{L}^{p}(\mathbf{R}^{n})\) 中的像（若 \(p \neq +\infty\) ）稠密。

对 \(p = +\infty\) ，第一个断言是直接的。现设 \(p \in [1, +\infty[\) 。设 m 是满足 \(n + 1 < mp\) 的整数。对每个 \(\varphi \in \mathcal{S}(\mathbf{R}^{n})\) 与 \(x \in \mathbf{R}^{n}\) ，我们有

\[
\| x \| ^ {n + 1} | \varphi (x) | ^ {p} \leqslant \widetilde {q} _ {0, m} (\varphi) ^ {p}
\]

因而依据 IV, p. 199 的命题 3，\(\varphi \in \mathcal{L}^p (\mathbf{R}^n)\) 。此外，我们得到

\[
\mathrm{N} _ {p} (\varphi) \leqslant a _ {n} ^ {1 / p} \widetilde {q} _ {0, 0} (\varphi) + b _ {n} \widetilde {q} _ {0, m} (\varphi),
\]

其中

\[
a _ {n} = \int_ {\| x \| \leqslant 1} d \mu (x), \quad b _ {n} = \int_ {\| x \| \geqslant 1} \frac {1}{\| x \| ^ {n + 1}} d \mu (x),
\]

因此 \(\mathcal{S}(\mathbf{R}^n)\) 到 \(\mathcal{L}^p (\mathbf{R}^n)\) 中的单射是连续的。

由于空间 \(\mathcal{D}(\mathbf{R}^n)\) 含于 \(\mathcal{S}(\mathbf{R}^n)\) ，IV, p. 202 的命题 4 蕴含 \(\mathcal{S}(\mathbf{R}^n)\) 在 \(\mathrm{L}^p (\mathbf{R}^n)\) 中稠密（若 \(p\neq +\infty\) ）。

\subsection{多项式增长的函数}

定义 3. --- 函数 \(f\colon\mathbf{R}^{n}\to\mathbf{C}\) 称为具有多项式增长的，如果存在整数 \(k\geqslant1\) ，使得由 \(x\mapsto(1+\|x\|)^{-k}f(x)\) 定义的映射在 \(\mathbf{R}^{n}\) 上有界。

每个具有多项式增长的函数都是局部有界的。\(\mathbf{R}^{n}\) 上的每个多项式函数都具有多项式增长。

命题 14. --- 设 \(f \in \mathcal{C}^{\infty}(\mathbf{R}^{n})\) 。假设对每个 \(\alpha\) （\(\mathbf{N}^{n}\) 中的），函数 \(\partial^{\alpha}f\) 都具有多项式增长。那么从空间 \(\mathcal{S}(\mathbf{R}^{n})\) 到自身的、由 \(\varphi \mapsto f\varphi\) 所定义的线性映射是连续的。

对每个 \(\varphi\) （在 \(\mathcal{S}(\mathbf{R}^{n})\) 中），函数 \(f\varphi\) 属于 \(\mathcal{C}^{\infty}(\mathbf{R}^{n})\) 。由假设，对每个 \(\alpha \in \mathbf{N}^{n}\) ，存在整数 \(k_{\alpha} \geqslant 0\) 与实数 \(\mathrm{C}_{\alpha}\) ，使得 \(|\partial^{\alpha}f(x)| \leqslant \mathrm{C}_{\alpha}(1 + \|x\|)^{k_{\alpha}}\) 对所有 \(x\) （在 \(\mathbf{R}^{n}\) 中）成立。设 \(\varphi \in \mathcal{S}(\mathbf{R}^{n})\) 。设 \(\alpha \in \mathbf{N}^{n}\) 与 \(k \in \mathbf{N}\) 。依据 IV, p. 196 的命题 1，可得

\[
\begin{array}{l} \widetilde {q} _ {\alpha , k} (f \varphi) \leqslant \sum_ {0 \leqslant \beta \leqslant \alpha} \binom {\alpha} {\beta} \sup _ {x \in \mathbf {R} ^ {n}} \| x \| ^ {k} | \partial^ {\beta} f (x)   \partial^ {\alpha - \beta} \varphi (x) | \\ \leqslant \sum_ {0 \leqslant \beta \leqslant \alpha} \binom {\alpha} {\beta} C _ {\beta} \sup _ {x \in \mathbf {R} ^ {n}} \| x \| ^ {k} (1 + \| x \|) ^ {k _ {\beta}} | \partial^ {\alpha - \beta} \varphi (x) |. \end{array}
\]

设 \(\beta \in \mathbf{N}^n\) 满足 \(0 \leqslant \beta \leqslant \alpha\) 。对每个 \(x \in \mathbf{R}^n\) ，我们有

\[
\begin{array}{c} \| x \| ^ {k} (1 + \| x \|) ^ {k _ {\beta}} | \partial^ {\alpha - \beta} \varphi (x) | \leqslant \sup _ {\| x \| \leqslant 1} 2 ^ {k _ {\beta}} | \partial^ {\alpha - \beta} \varphi (x) | \\ + \sup _ {x \in \mathbf {R} ^ {n}} 2 ^ {k _ {\beta}} \| x \| ^ {k + k _ {\beta}} | \partial^ {\alpha - \beta} \varphi (x) | \end{array}
\]

于是最终得到

\[
\widetilde {q} _ {\alpha , k} (f \varphi) \leqslant \sum_ {0 \leqslant \beta \leqslant \alpha} \binom {\alpha} {\beta} 2 ^ {k _ {\beta}} \mathrm{C} _ {\beta} \Bigl (\widetilde {q} _ {\alpha - \beta , 0} (\varphi) + \widetilde {q} _ {\alpha - \beta , k + k _ {\beta}} (\varphi) \Bigr),
\]

由此即得命题。

\subsection{缓增广义函数}

定义 4. --- 我们把 \(\mathbf{R}^{n}\) 上的缓增广义函数空间称为赋以有界收敛拓扑的 \(\mathcal{S}(\mathbf{R})\) 的对偶空间。记作 \(\mathcal{S}'(\mathbf{R}^{n})\) 。

由于 \(\mathcal{S}(\mathbf{R}^n)\) 是有界型的，空间 \(\mathcal{S}'(\mathbf{R}^n)\) 是完备的（TVS, III, p. 24，推论 1）。由于 \(\mathcal{S}(\mathbf{R}^n)\) 是蒙泰尔空间，\(\mathcal{S}'(\mathbf{R}^n)\) 亦然（EVT, IV, p. 19，命题 9）。于是空间 \(\mathcal{S}'(\mathbf{R}^n)\) 是自反的（EVT, IV, p. 16，定理 2）。

引理 9. --- 从 \(\mathcal{S}(\mathbf{R}^{n})\) 到 \(\mathbf{C}\) 中的线性映射 f 是缓增广义函数当且仅当：对每个族 \((\mathrm{M}_{\alpha,k})_{(\alpha,k)\in\mathbf{N}^{n}\times\mathbf{N}}\) （\(\mathbf{R}_{+}\) 中的），线性形式 f 在由这样的函数 \(\varphi\in\mathcal{S}(\mathbf{R}^{n})\) ——满足 \(\widetilde{q}_{\alpha,k}(\varphi)\leqslant\mathrm{M}_{\alpha,k}\) （对所有 \((\alpha,k)\in\mathbf{N}^{n}\times\mathbf{N}\) ）——所成之集上有界。

事实上，空间 \(\mathcal{S}(\mathbf{R}^{n})\) 是有界型的（EVT, III, p. 12，命题 2），且 \(\mathcal{S}(\mathbf{R}^{n})\) 的每个有界子集都含于陈述中所描述的有界集之一。

引理 10. --- 设 F 是 \(\mathcal{S}'(\mathbf{R}^n)\) 上具有可数基或包含一个单有界集合的滤子。那么 F 收敛于某个缓增广义函数当且仅当 \(\langle f, \varphi \rangle\) 对每个施瓦茨函数 \(\varphi\) 依 F 收敛。

特别地，缓增广义函数序列 \((f_{m})_{m\in\mathbf{N}}\) 收敛于缓增广义函数 f 当且仅当 \(\langle f_{m},\varphi\rangle\to\langle f,\varphi\rangle\) 对所有 \(\varphi\in\mathcal{S}(\mathbf{R}^{n})\) 成立。

由于 \(\mathcal{S}'(\mathbf{R}^n)\) 是弗雷歇空间，故为桶型空间（EVT, III, p. 25，命题 2 的推论），引理由巴拿赫–斯坦豪斯定理得出（EVT, III, p. 26，定理 1 的推论 3）。

典范单射 \(j\) ——\(\mathcal{D}(\mathbf{R}^{n})\) 到 \(\mathcal{S}(\mathbf{R}^{n})\) 中的——是连续的且其像稠密（IV, p. 212，引理 12），因而 \(j\) 的转置——即缓增广义函数到子空间 \(\mathcal{D}(\mathbf{R}^{n})\) 的限制映射——是从 \(\mathcal{S}'(\mathbf{R}^{n})\) 到 \(\mathcal{D}'(\mathbf{R}^{n})\) 中的单连续线性映射。我们将通过此映射把 \(\mathcal{S}'(\mathbf{R}^{n})\) 与 \(\mathcal{D}'(\mathbf{R}^{n})\) 的一个子空间等同起来。

设 \(\alpha \in \mathbf{N}^n\) 。线性映射 \(\varphi \mapsto \partial^\alpha \varphi\) ——从 \(\mathcal{S}(\mathbf{R}^n)\) 到 \(\mathcal{S}(\mathbf{R}^n)\) 中——是连续的。因此它的转置是从 \(\mathcal{S}'(\mathbf{R}^n)\) 到 \(\mathcal{S}'(\mathbf{R}^n)\) 中的连续线性映射（EVT, IV, p. 6，推论，\(b\) ））。我们用 \(\partial^\alpha\) 表示连续线性映射 \((-1)^{|\alpha|}{}^{t}\partial^\alpha\) ——从 \(\mathcal{S}'(\mathbf{R}^n)\) 到 \(\mathcal{S}'(\mathbf{R}^n)\) 中。该定义与 IV, p. 208 中关于广义函数的定义 2 相容。

设 \(h \in \mathcal{C}^{\infty}(\mathbf{R}^{n})\) 是 \(\partial^{\alpha}h\) （对每个 \(\alpha \in \mathbf{N}^{n}\) ）都具有多项式增长的函数。连续线性映射 \(\varphi \mapsto h\varphi\) （IV, p. 214，命题 14）的转置是 \(\mathcal{S}'(\mathbf{R}^{n})\) 上的连续线性映射，记作 \(f \mapsto hf\) 。